\chapter{Change-Log}
\mpitermtitleindex{change-log}
\label{sec:change}
\label{chap:change}

Annex~\ref{sec:change:newest} summarizes changes from the previous
version of the
\MPI/ standard to the version presented by this document.
Only significant changes (i.e., clarifications and new features) 
that might either require implementation effort in the 
\MPI/ libraries or change the understanding of \MPI/ from a 
user's perspective are presented.
Editorial modifications, formatting, typo corrections and
minor clarifications are not shown. 
If not otherwise noted, the section and page references refer to the
locations of the change or new functionality in this
version of the standard. 
Changes in Annexes~\ref{sec:change:recent}--\ref{sec:changes:last}
were already introduced in the corresponding sections
in previous versions of this standard. 

%%
%% Use raggedright for these entries because they often include long routine
%% names and other items that make right justification hard.
%%
\section{Changes from Version 5.0 to Version 5.1}
\label{subsec:50to51}
%% Keep the label sec:change:newest with the most recent changes
\label{sec:change:newest}

\begin{enumerate}
% 01.--- MPI-5.0-issue #20 PR665
\item
    Chapter~\ref{chap:ft}, Sections~\ref{sec:terms-errorhandling}, \ref{subsec:inquiry-inquiry}, \ref{sec:errorhandler}, and \ref{sec:ei-error-classes},
    on pages~\pageref{chap:ft}, \pageref{sec:terms-errorhandling}, \pageref{subsec:inquiry-inquiry}, \pageref{sec:errorhandler}, and \pageref{sec:ei-error-classes}.
    \newline
    Added operations, procedures, constants, and semantics to handle \MPI/ process failures.
\end{enumerate}

\subsection{Fixes to Errata in Previous Versions of \texorpdfstring{\MPI/}{MPI}}
\label{sec:changes-50-51-errata}

\begin{enumerate}
% MPI-5.1 issue #1001 PR #1127
\item
  Section~\ref{subsec:abi-opaque-handles} on page~\pageref{subsec:abi-opaque-handles}, and \MPIVDOTO/ Section 20.3.2 on page 846.
  \newline
  Added missing handle type definitions for opaque handle types of the \MPI/ tool information interface.

\end{enumerate}

\subsection{Changes in \texorpdfstring{\MPIVDOTI/}{MPI-5.1}}
\label{subsec:changes-50-51}


\section{Changes from Version 4.1 to Version 5.0}
\label{subsec:41to50}
\label{sec:change:recent}
%% The label sec:change:recent should be moved to the section for the
%% changes the version immeadiately prior to this one - e.g.,
%% for MPI 5.1, following MPI 5.0, this label is attached to
%% the section on changes from MPI 4.1 to MPI 5.0. When the next
%% version of MPI is released (probably MPI 5.2), move this label to
%% the section covering changes between MPI 5.0 and MPI 5.1.
%% Keep the label sec:change:newest with the most recent changes

\subsection{Fixes to Errata in Previous Versions of \texorpdfstring{\MPI/}{MPI}}
\label{sec:changes-41-50-errata}

\begin{enumerate}
% 01.--- MPI-5.0-issue #889 PR1034 MPI-5.0-p39, 4.1-page 39
\item
  Section~\ref{subsec:pt2pt-status} on
  page~\pageref{subsec:pt2pt-status}, and \MPIIVDOTI/ Section~3.2.5 on page~39.
  \newline
  Mark \mpiarg{status} argument as \code{const} in \mpifunc{MPI\_STATUS\_GET\_SOURCE},
  \mpifunc{MPI\_STATUS\_GET\_TAG}, and \mpifunc{MPI\_STATUS\_GET\_ERROR}.
  
% 02.--- MPI-5.0-issue #810 PR946 MPI-5.0-p484, 4.1-page 484
\item
  Sections~\ref{sec:inquiry-startup} and~\ref{subsec:mem_alloc_kinds} on
  pages~\pageref{sec:inquiry-startup} and~\pageref{subsec:mem_alloc_kinds}, and
  \MPIIVDOTI/ Sections~11.2.1 and~11.4.3 on pages 484 and 508.
  \newline
  Clarified that, in the \wpm{}, the value of the \infokey{mpi\_memory\_alloc\_kinds}
  info key in \mpiconst{MPI\_INFO\_ENV} indicates the requested memory allocation kinds.
  The provided memory allocation kinds can be queried through the value of the \infokey{mpi\_memory\_alloc\_kinds}
  info key in the info object returned by a call to
  \mpifunc{MPI\_COMM\_GET\_INFO} on \mpiconst{MPI\_COMM\_WORLD} or \mpiconst{MPI\_COMM\_SELF}.

% 03.--- MPI-5.0-issue #773 PR986 MPI-5.0-p510, 4.1-page 508
\item
  Section~\ref{subsec:mem_alloc_kinds} on
  page~\pageref{subsec:mem_alloc_kinds}, and
  \MPIIVDOTI/ Section~11.4.3 on page 508.
  \newline
  Clarified the requirements of the \infokey{mpi\_memory\_alloc\_kinds}
  info key's value when defaulted.

% 04.--- MPI_5.0-issue #760 PR904 MPI-5.0-p580, 4.1-page 580
\item 
  Section~\ref{sec:1sided-req} page~\pageref{sec:1sided-req} and \MPIIVDOTI/~Section~12.3.5 on page~580.
  \newline
  The text contained contradicting statements about the use of \mpifunc{MPI\_REQUEST\_FREE} on \RMA/ requests.
  A statement that suggested the use of \mpifunc{MPI\_REQUEST\_FREE} to release \RMA/ requests after the window has been flushed was removed.

\end{enumerate}

\subsection{Changes in \texorpdfstring{\MPIVDOTO/}{MPI-5.0}}
\label{subsec:changes-41-50}

\begin{enumerate}

% 01.--- MPI-5.0-issue #699 PR963 MPI-5.0-p228
\item Fixed-size Fortran logical datatypes were added to \sectionref{subsec:pt2pt-messagedata},
      \sectionref{coll-predefined-op}, Table \ref{table:io:optional-extsizes},
      \sectionref{f90:types}, and Appendix \sectionref{sec:lang}.

% 02.--- MPI-5.0-issue #811 PR938, MPI-5.0-p314, MPI-4.1 page 318
\item
  Section~\ref{subsec:context-grpconst} on page~\pageref{subsec:context-grpconst}.
  \newline
  A restriction on the provenance of the process set name supplied to \mpifunc{MPI\_GROUP\_FROM\_SESSION\_PSET} was removed.

% 03.--- MPI-5.0-issue #751 PR977 MPI-5.0-p853
\item { \raggedright Chapter~\ref{chap:abi} was added, which defines new functions,
        \mpifunc{MPI\_ABI\_GET\_VERSION} and \mpifunc{MPI\_ABI\_GET\_INFO},
        the layout of the \MPI/ status object,
        and the type of \MPI/ object handles (e.g. \mpictype{MPI\_Comm}).
      Extensive changes and additions were made in Section~\ref{subsec:annexa-const}
        to define literal values for many \MPI/ constants.
      In Chapter~\ref{chap:info}, \mpiconst{MPI\_MAX\_INFO\_KEY} was
          modified by one to be consistent with other constants.
      \mpifunc{MPI\_ABI\_GET\_VERSION} and \mpifunc{MPI\_ABI\_GET\_INFO} were added to
          Table~\ref{table:dynamic:preinit_functions}. }

% 04.--- MPI-5.0-issue #654 PR978 MPI-5.0-p858
\item {\sloppy%%ALLOWLATEX%%
  \sectionref{subsec:fortran-type-registration} was added to define new functions,
        \mpifunc{MPI\_ABI\_SET\_FORTRAN\_INFO}, \mpifunc{MPI\_ABI\_GET\_FORTRAN\_INFO},
        and \mpifuncshort{MPI\_ABI\_SET\_FORTRAN\_BOOLEANS}, along with the associated info keys
        and an error code \error{MPI\_ERR\_ABI},
        in order to expose Fortran environment settings to the \MPI/ library.
      \sectionref{subsec:fortran-abi-names} and \sectionref{subsec:fortran-abi-status}
        were added to define basic properties of the \MPI/ Fortran ABI.
      \sectionref{sec:abi-handleserialization} was added to obviate the need for \mpictype{MPI\_Fint}.

}

\end{enumerate}

\section{Changes from Version 4.0 to Version 4.1}
\label{subsec:40to41}

\subsection{Fixes to Errata in Previous Versions of \texorpdfstring{\MPI/}{MPI}}
\label{sec:changes-40-41-errata}

\begin{enumerate}

% 01.--- MPI-4.1-issue #676 PR825 4.0-page 13
\item
  Sections~\ref{subsec:operations}, \ref{sec:pt2pt-semantics}, \ref{sec:pt2pt-nonblock}, \ref{sec:pt2pt-probe}, \ref{sec:nbcoll},
\ref{sec:context:mpilibsupport}, \ref{sec:context:contexts}, \ref{subsec:context-intracomconst},
\ref{sec:dynamic:threads:clarifications}, \ref{sec:io-filecntl-preallocate}, \ref{sec:io-semantics-nb-collective},
\ref{sec:f90-problems:code-movements}, \ref{sec:f90-problems:comparison-with-C},
  on pages~\pageref{subsec:operations}, \pageref{sec:pt2pt-semantics}, \pageref{sec:pt2pt-nonblock}, \pageref{sec:pt2pt-probe},
\pageref{sec:nbcoll}, \pageref{sec:context:mpilibsupport}, \pageref{sec:context:contexts}, \pageref{subsec:context-intracomconst},
\pageref{sec:dynamic:threads:clarifications}, \pageref{sec:io-filecntl-preallocate}, \pageref{sec:io-semantics-nb-collective},
\pageref{sec:f90-problems:code-movements}, \pageref{sec:f90-problems:comparison-with-C}
  and \MPIIVDOTO/ Sections~2.4.1, 3.5, 3.7, 3.8, 6.12, 7.1.2, 7.2.2, 7.4.2, 11.6.2, 14.2.5, 14.6.5, 19.1.17, 19.1.20
  on pages~13, 54, 60, 84, 250, 312, 314, 327, 518, 650, 713, 826, and 834.
  \newline
  The term \mpitermni{pending}\mpitermindex{pending (communication) operation}\mpitermindex{pending I/O operation}\mpitermindex{communication!pending}\mpitermindex{operation!pending}
  communication or I/O operation is defined as \mpitermni{in the active operation state}\mpitermindex{active!operation state}.
  If the phrase \mpitermni{pending communication operation} in \MPIIDOTI/ to \MPIIVDOTO/ additionally includes
  \mpitermni{decoupled \MPI/ activities}\mpitermindex{decoupled MPI activities@decoupled \MPI/ activities}, then this has
been added explicitly.
  If this phrase had a different meaning, it was replaced accordingly, see
  item~\ref{item:changes:issue710-PR823} in this list and  % (outcome of MPI\_COMM\_DISCONNECT)
  item~\ref{item:changes:issue543-PR791} in \sectionref{subsec:changes-40-41}.  % (definition of MPI\_COMM\_FREE was clarified)}

% 02.--- MPI-4.1-issue #705 and #657 fixed by PR822 4.0-page 20
\item
  Sections~\ref{subsec:named-constants} and~\ref{subsec:annexa-const} on pages~\pageref{subsec:named-constants} and~\pageref{subsec:annexa-const}, and
  \MPIIVDOTO/ Sections~2.5.4 and~A.1.1 on pages~20 and~857.
  \newline
  The implementation of named \MPI/ \mpiterm{constants} in C and Fortran and implied usage restrictions were clarified.

% 03.--- MPI-4.1-issue #658 PR776 4.0-page 20
\item
  Section~\ref{subsec:named-constants} on page~\pageref{subsec:named-constants}, and \MPIIVDOTO/ Section~2.5.4 on page~20.
  \newline
  Add \mpiconst{MPI\_MAX\_PSET\_NAME\_LEN} and \mpiconst{MPI\_MAX\_STRINGTAG\_LEN} to list of named constants.

% 04.--- MPI-4.1-issue #710 PR823 4.0-pages 94, 494 and 546
\item \label{item:changes:issue710-PR823}
  Sections~\ref{sec:pt2pt-persistent}, \ref{sec:inquiry-finalize}, \ref{sec:session-init-finalize} and~\ref{subsec:disconnect} 
  on pages~\pageref{sec:pt2pt-persistent}, \pageref{sec:inquiry-finalize}, \pageref{sec:session-init-finalize} and~\pageref{subsec:disconnect}
  and \MPIIVDOTO/ Sections~3.9, 11.2.2, 11.3.1 and~11.10.4 on pages~94, 494, 501 and~546.
  \newline
  The requirements for calling \mpifunc{MPI\_FINALIZE}, \mpifunc{MPI\_SESSION\_FINALIZE}, and \mpifunc{MPI\_COMM\_DISCONNECT} 
  and the outcome of \mpifunc{MPI\_COMM\_DISCONNECT}, especially for related
  inactive persistent request handles, were clarified.
  
% 05.--- MPI-4.1-issue #589 PR726 4.0-page 115
\item
  Section~\ref{sec:part-communication-examples} on page~\pageref{sec:part-communication-examples}, and
  \MPIIVDOTO/ Section~4.3.3. on page~115.
  Example~\ref{partexample4} on page~\pageref{partexample4}, and
  \MPIIVDOTO/ Example~4.4 on page 115. 
  \newline
  Fixed and simplified erroneous \MPIIVDOTO/ Example 4.4. The example could deadlock due to
  incorrect use of the \mpicode{flag} variable in multiple \MPI/ test procedure calls or
  thread concurrent access.
  The example was also simplified by removing unnecessary code and updated according to current
  best practice in OpenMP.

% 06.--- MPI 4.1-issue #509 PR607 4.0-page 122
\item 
Section~\ref{subsec:part-channel-examples} on page~\pageref{subsec:part-channel-examples}, \MPIIV/ Section~4.3.1 on page 112.
\newline
Example~\ref{part-example2} was corrected.

% 07.--- MPI-4.1-issue #710 PR823 4.0-pages 150, 488, 494 and 501
\item
  Section~\ref{subsec:pt2pt-comfree} on page~\pageref{subsec:pt2pt-comfree}
  and \MPIIVDOTO/ Section~5.1.9 on page~150.
\newline
The relationship between \mpifunc{MPI\_TYPE\_COMMIT} and initialization/finalization of a session (or the \wpm{})
with \mpifunc{MPI\_INIT}, \mpifunc{MPI\_INIT\_THREAD}, \mpifunc{MPI\_FINALIZE},
\mpifunc{MPI\_SESSION\_INIT}, and \mpifunc{MPI\_SESSION\_FINALIZE}
was clarified.

% 08--- MPI-4.1-issue #511 PR629 4.0-pages 343 and 360
\item
	Sections~\ref{subsec:context-intracomconst} on page~\pageref{subsec:context-intracomconst} and~\ref{subsec:context-intercomm} on page~\pageref{subsec:context-intercomm}, and \MPIIVDOTO/ Section 7.4.2 on page~343 and Section 7.6.2 on page~360.
  \newline
  Use of the \mpiarg{errhandler} argument to \mpifunc{MPI\_COMM\_CREATE\_FROM\_GROUP} and \mpifunc{MPI\_INTERCOMM\_CREATE\_FROM\_GROUPS} is clarified.
  The error handler invoked when an error is encountered during invocation of these functions is also clarified.

% 09.--- MPI 4.1-issue #690 PR830 4.0-page 328
\item
  Section~\ref{subsec:context-intracomconst} on page~\pageref{subsec:context-intracomconst}, \MPIIII/ Section~6.4.2 on page 237,
  \MPIIIIDOTI/ Section~6.4.2 on page 237, and \MPIIV/ Section~7.4.2 on page 327. 
  \newline
  The description of \mpifunc{MPI\_COMM\_DUP} now clarifies that error handlers are also copied in
  the output communicator produced when this procedure is called.

% 10.--- MPI 4.1-issue #687 PR#827 4.0-page 369
\item
  Section~\ref{subsec:context-cachecomm} on page~\pageref{subsec:context-cachecomm}, \MPIIII/ Section~6.7.2 on page 267, \MPIIIIDOTI/ Section~6.7.2 on page 267, and \MPIIV/ Section~7.7.2 on page 366.
  \newline
  \mpifunc{MPI\_COMM\_DISCONNECT} was explicitlly stated as a procedure that deletes attributes.
  
% 11.--- MPI-4.1-issue #556 PR690 4.0-page 403 (406-407)
\item
  Section~\ref{subsec:topol-inquiry} on page~\pageref{subsec:topol-inquiry} and \MPIIVDOTO/ Section 8.5.5 on page 403. 
  \newline
  The unintended change in the specification of argument \mpiarg{coords} in \mpifunc{MPI\_CART\_COORDS} in \MPIIVDOTO/ is reverted to the original meaning in \MPIIDOTI/ to \MPIIIIDOTI/. It is clarified that the outcome of               \mpifunc{MPI\_CART\_GET} and \mpifunc{MPI\_CART\_COORDS} is unspecified for the case that \mpiarg{maxdims} is less than \mpiarg{ndims}.

% 12.--- MPI-4.1-issue #588 PR644 4.0-page 459
\item
  Section~\ref{sec:errorhandler} on page~\pageref{sec:errorhandler} and \MPIIVDOTO/ Section 9.3 on page 458.
  \newline
  The fallback error-handler for the \spm{} was clarified.

% 13.--- MPI-4.1-issue #545 PR657 4.0-page 473
\item
  Section~\ref{sec:ei-error} on page~\pageref{sec:ei-error} and \MPIIVDOTO/ Section~9.5 on page~473.
  \newline
  It was clarified that \mpiconst{MPI\_LASTUSEDCODE} is only available in the \wpm{}.

% 14.--- MPI-4.1-issue #434 PR618 4.0-page 487
\item
  Sections~\ref{sec:dynamic:introduction} and~\ref{sec:model} on pages~\pageref{sec:dynamic:introduction} and~\pageref{sec:model},
  and \MPIIVDOTO/ Sections~11.1 and~11.7 on pages~487 and~521.
  \newline
  It was clarified that the usage of the Dynamic Process Model requires the \wpm{} to be initialized.

% 15.--- MPI-4.1-issue #499 PR729 4.0-page 601
\item
  Section~\ref{sec:1sided-sync:pscw} on page~\pageref{sec:1sided-sync:pscw}  and \MPIIV/ Section~12.5.2 on page~598.
  \newline
  The definition of \mpifunc{MPI\_WIN\_TEST} was clarified.

% 16.--- MPI 4.1-issue #634 PR747 4.0-page 607
\item \mpitermindex{progress!load/store access synchronization}
Section~\ref{sec:1sided-flush} on page~\pageref{sec:1sided-flush}, \MPIIII/ Section~11.5.4 on page 449, \MPIIIIDOTI/ Section~11.5.4 on page 448, and \MPIIV/ Section~12.5.4 on page 605.
\newline
The description of \mpifunc{MPI\_WIN\_SYNC} was clarified to include its use for ordering load/store accesses to shared memory.
A statement was added to highlight that a call to \mpifunc{MPI\_WIN\_SYNC} does not complete pending \RMA/ operations and
that a call to \mpifunc{MPI\_WIN\_SYNC} does not guarantee any progress of \MPI/ operations.

% 17.--- MPI-4.1-issue #517 PR620 4.0-page 640
\item 
  Section~\ref{sec:ei-status} on page~\pageref{sec:ei-status}, and \MPIIV/ Section~13.3 on page 640.
\newline
  Large count interface of \mpifunc{MPI\_STATUS\_SET\_ELEMENTS} had been missing and was added.

% 18.--- MPI_4.0-issue #443 PR754
\item Sections~\ref{sec:mpit:cvar} and~\ref{sec:mpit:pvar} on pages~\pageref{sec:mpit:cvar} and~\pageref{sec:mpit:pvar}, \MPIIII/ Sections~14.3.6 and 14.3.7 on pages~567 and 573, \MPIIIIDOTI/ Sections~14.3.6 and 14.3.7 on pages~573 and 580, and \MPIIV/ Sections~15.3.6 and 15.3.7 on pages~738 and 744.
\newline
The intent of handle arguments of the language independent definition of
\mpifunc{MPI\_T\_CVAR\_WRITE}, %
\mpifunc{MPI\_T\_PVAR\_HANDLE\_ALLOC},\flushline %
\mpifunc{MPI\_T\_PVAR\_HANDLE\_FREE}, %
\mpifunc{MPI\_T\_PVAR\_START}, %
\mpifunc{MPI\_T\_PVAR\_STOP}, %
\mpifunc{MPI\_T\_PVAR\_WRITE}, and %
\mpifunc{MPI\_T\_PVAR\_RESET}%
were marked as \INOUT\ in accordance with the special rule for handles described in Section~\ref{terms:procedure-specification}.

\end{enumerate}

\subsection{Changes in \texorpdfstring{\MPIIVDOTI/}{MPI-4.1}}
\label{subsec:changes-40-41}

\begin{enumerate}


% 01.--- MPI-4.1-issue #562 PR694, MPI-4.0 page 6 (new section will start there)
\item
  Section~\ref{sec:side-documents} on page~\pageref{sec:side-documents}.
  \newline
  Introduced the concept of side documents.

% 02.--- MPI-4.1-issue #669 PR788, MPI-4.0 pages 14, 15, and 17
\item
  Sections~\ref{subsec:operations} and~\ref{subsec:procedures} on pages~\pageref{subsec:operations} and~\pageref{subsec:procedures}.
  \newline
  Added the definition of the
  \mpiterm{enabled operation state}\mpitermindex{state!operation!enabled}\mpitermindex{operation!state!enabled}.

% 03.--- MPI-4.1-issue #565 PR727
\item
  Sections~\ref{sub:choice}, \ref{sec:fortran-binding-issues}, and \ref{sec:fortran}
  on pages~\pageref{sub:choice}, \pageref{sec:fortran-binding-issues}, and \pageref{sec:fortran}.
  \newline
  The \MPI/ standard now reflects that TS 29113\mpitermindex{TS 29113}
  was superceded by Fortran 2018\mpitermindex{Fortran 2018}.

% 04.--- MPI-4.1-issue #518 (PR623) PR832 MPI-4.0 on pages 23, 141, 147, 149, 153, 640, (786)
\item
    Sections~\ref{sec:deprecated}, \ref{subsec:pt2pt-addfunc}, \ref{subsec:pt2pt-extent}, \ref{subsec:pt2pt-true-extent}, \ref{subsec:pt2pt-datatypeuse}, \ref{sec:ei-status},
    and~\ref{sec:deprecated:since41}
    on pages~\pageref{sec:deprecated}, \pageref{subsec:pt2pt-addfunc}, \pageref{subsec:pt2pt-extent}, \pageref{subsec:pt2pt-true-extent}, \pageref{subsec:pt2pt-datatypeuse}, \pageref{sec:ei-status},
    and~\pageref{sec:deprecated:since41}.
    \newline
    \mpifunc{MPI\_TYPE\_SIZE\_X}, \mpifunc{MPI\_TYPE\_GET\_EXTENT\_X},
    \newline
    \mpifunc{MPI\_TYPE\_GET\_TRUE\_EXTENT\_X}, \mpifunc{MPI\_GET\_ELEMENTS\_X}, and
    \newline
    \mpifunc{MPI\_STATUS\_SET\_ELEMENTS\_X}
    were deprecated and may be removed in a future version of the \MPI/ specification.

% 05.--- MPI 4.1-issue #57 PR790
\item
  Sections~\ref{sec:deprecated}, \ref{sec:predef-comms} and~\ref{subsec:inquiry-inquiry} on pages~\pageref{sec:deprecated}, \pageref{sec:predef-comms} and~\pageref{subsec:inquiry-inquiry}.
  \newline
  \mpiconst{MPI\_HOST} has been deprecated, and a mention to host process has been removed.

% 06.--- MPI-4.1-issue #561 PR722
\item
  Sections~\ref{sec:deprecated}, \ref{f90:overview}, \ref{f90:basic}, \ref{sec:deprecated:since41}
  on pages~\pageref{sec:deprecated}, \pageref{f90:overview}, \pageref{f90:basic}, \pageref{sec:deprecated:since41}.
  \newline
  Deprecated the use of \code{mpif.h}.

% 07.--- MPI-4.1-issue #646 PR771
\item
  Section~\ref{sec:macros} on page~\pageref{sec:macros}.
  \newline
  Removed the functions \mpifunc{MPI\_WTIME}, \mpifuncskip{PMPI\_WTIME},
  \mpifunc{MPI\_WTICK}, and \mpifuncskip{PMPI\_WTICK} from the list of functions
  that may be implemented as a macro.

% 08.--- MPI-4.1-issue #647 PR770
\item
  Section~\ref{sec:macros} on page~\pageref{sec:macros} and
  Section~\ref{sec:misc-handleconvert} on page~\pageref{sec:misc-handleconvert}.
  \newline
  Removed the ability to implement \mpi/ handle conversion functions as a macro.

% 09.--- MPI-4.1-issue #493 PR753 and MPI-4.1-issue #468 PR667, MPI-4.0 page 83
\item
  Sections~\ref{subsec:terms:progress}, \ref{subsec:pt2pt-teststatus}, \ref{sec:1sided-pscw},
  and Example~\ref{exa:sync-shared:deadlock}
  on pages~\pageref{subsec:terms:progress}, \pageref{subsec:pt2pt-teststatus}, \pageref{sec:1sided-pscw},
  and \pageref{exa:sync-shared:deadlock}.
  \newline
  The \mpitermni{progress rules}\mpitermindex{progress}\mpitermindex{progress!point-to-point communication}
  \mpitermindex{semantics!point-to-point communication!progress} were clarified in general
  and for \mpifunc{MPI\_REQUEST\_GET\_STATUS} and \mpifunc{MPI\_WIN\_TEST}.
  The terms \mpitermni{strong}\mpitermindex{strong progress}\mpitermindex{progress!strong}
  and \mpiterm{weak progress}\mpitermindex{progress!weak} were introduced.
  An example showing restrictions on the use of \MPI/ shared memory for
  synchronizing\mpitermindex{shared memory!loads/stores for synchronizing} purposes was introduced.

% 10.--- MPI 4.1-issue #645 PR767
\item
  Sections~\ref{subsec:pt2pt-status} and~\ref{sec:ei-status} on pages~\pageref{subsec:pt2pt-status} and~\pageref{sec:ei-status}.
  \newline
  Added procedures \mpifunc{MPI\_STATUS\_GET\_SOURCE}, \mpifunc{MPI\_STATUS\_GET\_TAG}, and \mpifunc{MPI\_STATUS\_GET\_ERROR} to query \MPI/ status fields
  and procedures \mpifunc{MPI\_STATUS\_SET\_SOURCE}, \mpifunc{MPI\_STATUS\_SET\_TAG}, and \mpifunc{MPI\_STATUS\_SET\_ERROR} to set \MPI/ status fields.
  Direct access to these fields remains available.

% 11.--- MPI 4.1-issue #586 PR736 and #685 PR808, Section 3.6, MPI-4.0 page 57
\item Section~\ref{sec:pt2pt-buffer} on page~\pageref{sec:pt2pt-buffer}.
  \newline
  % Common change-log entry for both issues and their PRs
  Automatic (unlimited) buffering is added, which can be enabled by using \mpiconst{MPI\_BUFFER\_AUTOMATIC} in any of the buffer attach procedures.
  New procedures \mpifunc{MPI\_COMM\_ATTACH\_BUFFER}, \mpifunc{MPI\_SESSION\_ATTACH\_BUFFER},
  \mpifunc{MPI\_COMM\_DETACH\_BUFFER} and \mpifunc{MPI\_SESSION\_DETACH\_BUFFER}
  % to attach and detach a buffer to a communicator or session, respectively, to be used for buffered sends
  % on that communicator or on communicators derived from groups derived from process sets from that session.
  have been\flushline % Needed to address overfull box
 added.
  The buffers attached with the existing functions \mpifunc{MPI\_BUFFER\_ATTACH} and \mpifunc{MPI\_BUFFER\_DETACH}
  now only apply to communicators that have no buffer attached at the communicator or session level.
  New procedures \mpifunc{MPI\_COMM\_FLUSH\_BUFFER}, \mpifunc{MPI\_SESSION\_FLUSH\_BUFFER}, and
  \mpifunc{MPI\_BUFFER\_FLUSH} were added as a combination of detach and attach, as well as the corresponding
  nonblocking variants \mpifunc{MPI\_COMM\_IFLUSH\_BUFFER}, \mpifunc{MPI\_SESSION\_IFLUSH\_BUFFER}, and
  \mpifunc{MPI\_BUFFER\_IFLUSH}.

% 12.--- MPI-4.1-issue #519 PR799
\item
 Subsection~\ref{subsec:pt2pt-teststatus} on page~\pageref{subsec:pt2pt-teststatus}
 \newline
 \begin{raggedright}
 Added new procedures \mpifunc{MPI\_REQUEST\_GET\_STATUS\_ANY}, \mpifunc{MPI\_REQUEST\_GET\_STATUS\_SOME},
 and \mpifunc{MPI\_REQUEST\_GET\_STATUS\_ALL} to \gb{}query the statuses of multiple requests without freeing them.
 \end{raggedright}

% 13.--- MPI-4.1-issue #520 PR798
\item
  Section~\ref{subsec:typeaccess} and~\ref{coll-minloc-maxloc} on pages~\pageref{subsec:typeaccess} and~\pageref{coll-minloc-maxloc}.
  \newline
  Added procedure \mpifunc{MPI\_TYPE\_GET\_VALUE\_INDEX} to query predefined
  datatype handles for pairs of value and index types to be usable in
  conjunction with \mpiconst{MPI\_MINLOC} and \mpiconst{MPI\_MAXLOC}.
  Added combiner \mpiconst{MPI\_COMBINER\_VALUE\_INDEX} for unnamed type handles
  returned by \mpifunc{MPI\_TYPE\_GET\_VALUE\_INDEX}.

% 14.--- MPI-4.1-issue #538 PR761, MPI-4.0 page 327
\item
  Section~\ref{subsec:context-intracomconst} on page~\pageref{subsec:context-intracomconst}.
  \newline
  \mpiconst{MPI\_COMM\_TYPE\_RESOURCE\_GUIDED} was added as a new possible value for the \mpiarg{split\_type}
  parameter of the \mpifunc{MPI\_COMM\_SPLIT\_TYPE} procedure, as well as a new info key
  \infokey{mpi\_pset\_name}.

% 15.--- MPI-4.1-issue #543 PR791, MPI-4.0 page 345
\item\label{item:changes:issue543-PR791}
  Section~\ref{subsec:context-intracomdest} on page~\pageref{subsec:context-intracomdest}.
  \newline
  The definition of \mpifunc{MPI\_COMM\_FREE} was clarified.

% 16.--- MPI-4.1-issue #83 PR666
\item
   Section~\ref{sec:comm-info} on page~\pageref{sec:comm-info}.
   \newline
   A new info key was added, namely
   \infokey{mpi\_assert\_strict\_persistent\_collective\_ordering}.

% 17.--- MPI-4.1-issue #580 PR714
\item
    Sections~\ref{sec:comm-info}, \ref{sec:inquiry-startup}, \ref{sec:model_sessions}, \ref{subsec:mem_alloc_kinds}, \ref{subsec:spawnkeys}, \ref{chap:one-side-2:win_create}, and~\ref{sec:io-info}
    on pages~\pageref{sec:comm-info}, \pageref{sec:inquiry-startup}, \pageref{sec:model_sessions}, \pageref{subsec:mem_alloc_kinds}, \pageref{subsec:spawnkeys}, \pageref{chap:one-side-2:win_create},~and \pageref{sec:io-info}.
    \newline
    Added the ability to request support for, query support of, and assert
    usage of memory allocation kinds via two new info keys,
    \infokey{mpi\_memory\_alloc\_kinds} and
    \infokey{mpi\_assert\_memory\_alloc\_kinds}.

% 18.--- MPI-4.1-issue #544 PR652 4.0-page 381 (383, 385-386)
\item
  Section~\ref{sec:ei-naming} on page~\pageref{sec:ei-naming}
  was amended to allow \mpiconst{MPI\_COMM\_NULL},
  \mpiconst{MPI\_DATATYPE\_NULL}, and \mpiconst{MPI\_WIN\_NULL} to be
  passed to \mpifunc{MPI\_COMM\_GET\_NAME},
  \mpifunc{MPI\_TYPE\_GET\_NAME} and
  \mpifunc{MPI\_WIN\_GET\_NAME}, respectively.

% 19.--- MPI 4.1-issue #154 PR762, MPI-4.0 page 453
\item
  Section~\ref{subsubsec:inquiry-hw-resource-info} on page~\pageref{subsubsec:inquiry-hw-resource-info}.
  \newline
  Added new procedure \mpifunc{MPI\_GET\_HW\_RESOURCE\_INFO}.

% 20.--- MPI-4.1-issue #525 PR633, MPI-4.0 page 469

\item
  Section~\ref{sec:ei-error-classes} on page~\pageref{sec:ei-error-classes}.
  \newline
  Added new error class
  \error{MPI\_ERR\_ERRHANDLER}.


% 21.--- MPI-4.1-issue #283 PR166, MPI-4.0 page 473
\item
 Section~\ref{sec:ei-error} on page~\pageref{sec:ei-error}.
 \newline
 Add procedures \mpifunc{MPI\_REMOVE\_ERROR\_CLASS}, \mpifunc{MPI\_REMOVE\_ERROR\_CODE},
 \mpifunc{MPI\_REMOVE\_ERROR\_STRING} to complement the procedures adding error classes, codes, and strings.

% 22.--- MPI-4.1-issue #23 PR708
\item
  Section~\ref{sec:one-side-2:intro} on page~\pageref{sec:one-side-2:intro}.
  \newline
  Relaxed the constraints on the windows for which shared memory can be queried using
  \mpifunc{MPI\_WIN\_SHARED\_QUERY} to allow windows with flavor \mpiconst{MPI\_WIN\_FLAVOR\_CREATE}
  and \mpiconst{MPI\_WIN\_FLAVOR\_ALLOCATE}.

% 23.--- MPI-4.1-issue #640 PR749, MPI-4.0 page 553 
\item
  Section~\ref{chap:one-side-2:win_create} on page~\pageref{chap:one-side-2:win_create}.
  \newline
  Added new info key \infokey{mpi\_accumulate\_granularity} to specify a desired
  synchronization granularity of accumulate operations.

% 24.--- MPI-4.1-issue #551 PR659, MPI-4.0 page 565

\item
  Section~\ref{sec:rma-win-free} on page~\pageref{sec:rma-win-free}.
  \newline
  Implementations may avoid synchronization of processes in \mpifunc{MPI\_WIN\_FREE}
  if the \infokey{no\_locks} info key is set to \infoval{true}.

% 25.--- MPI 4.1-issue #555 PR671
\item
  Section~\ref{sec:1sided-pscw} on page~\pageref{sec:1sided-pscw}.
  \newline
  In Example~\ref{ex:1sided-start-complete}, \mpifunc{MPI\_PUT} has been removed from the list
  of procedures that may delay their return waiting for the call to \mpifunc{MPI\_WIN\_POST} to occur at the target.
  \MPI/ \RMA/ communication procedures are generally not intended to delay their return waiting for synchronization procedure calls
  to occur at the target.

% 26.--- MPI-4.1-issue #553 PR663
\item
  Section~\ref{sec:1sided-semantics} on page~\pageref{sec:1sided-semantics} and
  Section~\ref{sec:onesided_examples} on page~\pageref{sec:onesided_examples}.
  \newline
  Clarified the use of \mpifunc{MPI\_WIN\_SYNC} for memory synchronization on shared memory.

% 27.--- MPI 4.1-issue #160 PR161 4.0-page 732
\item
  Section~\ref{sec:mpit:assoc} on page~\pageref{sec:mpit:assoc}.
  \newline
  The text specifying when entities of the \MPI/ Tool Information Interface can 
  be bound to objects during the object's lifetime was clarified.

% 28.--- MPI 4.1-issue #414 PR723
\item
    Section~\ref{sec:tools:mpit:droppedevents} on page~\pageref{sec:tools:mpit:droppedevents} and
    Section~\ref{sec:mpit:category_changed} on page~\pageref{sec:mpit:category_changed}.
    \newline
    Behavior specified when the count of dropped events or category changes overflow, respectively.

% 29.--- MPI-4.1-issue #535 PR658
\item
  Section~\ref{sec:misc-lang-interop:initialization} on page~\pageref{sec:misc-lang-interop:initialization}.
  \newline
  Language interoperability when using \mpifunc{MPI\_SESSION\_INIT} and \mpifunc{MPI\_SESSION\_FINALIZE} was clarified.

% 30.--- MPI-4.1-issue #641 PR752
\item
  Annex~\ref{sec:lang:semantics:summary} on page~\pageref{sec:lang:semantics:summary}.
  \newline
  The annex has been completed with the operation-related
  \MPI/ procedures for one-sided communication and some other rarely used scenarios.
  It is now integrated into the \MPI/ standard.

\end{enumerate}

\section{Changes from Version 3.1 to Version 4.0}
\label{subsec:31to40}

\subsection{Fixes to Errata in Previous Versions of \texorpdfstring{\MPI/}{MPI}}
\label{sec:changes-31-40-errata}

\begin{enumerate}

% 01.--- MPI-3.0-issue #153 p410
\item
\mpitermindex{neighborhood collective communication!periodic and dims 1 or 2@periodic and dims=1 or 2}
Sections~\ref{subsec:sparsecoll:allgather}, \ref{subsec:sparsecoll:alltoall} and~\ref{sec:topol-applic-example}
on pages~\pageref{subsec:sparsecoll:allgather}, \pageref{subsec:sparsecoll:alltoall} and~\pageref{sec:topol-applic-example},
and \MPIIIIDOTI/ Sections~7.6.1, 7.6.2 and 7.8 on pages 315, 318 and 329.
\newline
\mpifuncindex{MPI\_NEIGHBOR\_ALLTOALL}%
\mpifuncindex{MPI\_NEIGHBOR\_ALLTOALLV}%
\mpifuncindex{MPI\_NEIGHBOR\_ALLTOALLW}%
\mpiskipfunc{MPI\_NEIGHBOR\_ALLTOALL\{$|$V$|$W\}}
and
\mpifuncindex{MPI\_NEIGHBOR\_ALLGATHER}%
\mpifuncindex{MPI\_NEIGHBOR\_ALLGATHERV}%
\mpiskipfunc{MPI\_NEIGHBOR\_ALLGATHER\{$|$V\}}
for Cartesian virtual grids were clarified. 
An advice to implementors was added to illustrate a correct implementation
for the case of \mpicode{periods[d]=1} or \mpicode{.TRUE.} 
and \mpicode{dims[d]=1} or \mpicode{2} in a direction \mpicode{d}.

% 02.--- MPI-4.0-issue #298 p832.1
\item
  Section~\ref{sec:conversion:status} on
  page~\pageref{sec:conversion:status-f2f08-not-in-mpifh}, and
  \MPIIIIDOTI/ Section 17.2.5 on page 657 line 11.
  \newline
  Clarified that the \mpifunc{MPI\_STATUS\_F2F08} and
  \mpifunc{MPI\_STATUS\_F082F} routines and the declaration for
  \mpiftype{TYPE(MPI\_Status)} are not supposed to appear with
  \code{mpif.h}.

% 03.--- MPI-4.0-issue #306 p832.2

\item
  Sections~\ref{subsec:named-constants:F_STATUS_SIZE}, \ref{sec:conversion:status:C_F_constants},
  and~\ref{subsec:annexa-const:MPI_F_STATUS-constants} on
  pages~\pageref{subsec:named-constants:F_STATUS_SIZE},
  \pageref{sec:conversion:status:C_F_constants},
  and~\pageref{subsec:annexa-const:MPI_F_STATUS-constants}, and
  \MPIIIIDOTI/ Sections 2.5.4, 17.2.5, and A.1.1 on pages 15, 656, and
  669.
  \newline
  Define the C constants \mpiconst{MPI\_F\_STATUS\_SIZE}, \mpiconst{MPI\_F\_SOURCE},
  \mpiconst{MPI\_F\_TAG}, and \mpiconst{MPI\_F\_ERROR}.

% 04.--- MPI-4.0-issue #307 p832.3
\item
  Section~\ref{sec:conversion:status:status_f2f08} on
  page~\pageref{sec:conversion:status:status_f2f08}, and
  \MPIIIIDOTI/ Section 17.2.5 on page 658.
  \newline
  Added missing \code{const} to \IN\ parameters for
  \mpifunc{MPI\_STATUS\_F2F08} and \mpifunc{MPI\_STATUS\_F082F}.

\end{enumerate}

\subsection{Changes in \texorpdfstring{\MPIIVDOTO/}{MPI-4.0}}
\label{subsec:changes-31-40}

\begin{enumerate}

% 01.--- MPI-4.0-issue #405 p9
\item
  Sections~\ref{subsec:terms:naming}, \ref{sec:incompatible:since40}, and~\ref{sec:f90:linker-names}
  on pages~\pageref{subsec:terms:naming}, \pageref{sec:incompatible:since40}, and~\pageref{sec:f90:linker-names}.
  \newline
  The limit for the maximum length of \MPI/ identifiers was removed.

% 02.--- MPI-4.0-issue #96 p11
\item
\mpitermindex{semantics} Section~\ref{terms:semantic},
\ref{sec:pt2pt-modes}, \ref{subsec:pt2pt-commstart},
\ref{subsec:pt2pt-commend}, \ref{subsec:pt2pt-probe},
\ref{sec:matching-probe},
\ref{sec:coll-persistent},
\ref{sec:io-split-collective}, and
Annex~\ref{sec:lang:semantics:summary} on
pages~\pageref{terms:semantic}, \pageref{sec:pt2pt-modes},
\pageref{subsec:pt2pt-commstart}, \pageref{subsec:pt2pt-commend},
\pageref{subsec:pt2pt-probe}, \pageref{sec:matching-probe},
\pageref{sec:coll-persistent},
\pageref{sec:io-split-collective}, and
\pageref{sec:lang:semantics:summary}.
\newline
The semantic terms were updated.

% 03.--- MPI-4.0-issue #137 p11
\item Sections~\ref{subsec:count} and~\ref{sec:largecount} on pages~\pageref{subsec:count} and~\pageref{sec:largecount}, and throughout the entire document.
  \newline
  New large count functions \mpifuncskip{MPI\_\{\ldots\}\_c}
\MPIbindindexuse{MPI\_Send\_c}%
\MPIbindindexuse{MPI\_Recv\_c}%
\MPIbindindexuse{MPI\_Get\_count\_c}%
\MPIbindindexuse{MPI\_Sendrecv\_c}%
\MPIbindindexuse{MPI\_Sendrecv\_replace\_c}%
\MPIbindindexuse{MPI\_Bsend\_c}%
\MPIbindindexuse{MPI\_Ssend\_c}%
\MPIbindindexuse{MPI\_Rsend\_c}%
\MPIbindindexuse{MPI\_Buffer\_attach\_c}%
\MPIbindindexuse{MPI\_Buffer\_detach\_c}%
\MPIbindindexuse{MPI\_Isend\_c}%
\MPIbindindexuse{MPI\_Ibsend\_c}%
\MPIbindindexuse{MPI\_Issend\_c}%
\MPIbindindexuse{MPI\_Irsend\_c}%
\MPIbindindexuse{MPI\_Irecv\_c}%
\MPIbindindexuse{MPI\_Isendrecv\_c}%
\MPIbindindexuse{MPI\_Isendrecv\_replace\_c}%
\MPIbindindexuse{MPI\_Mrecv\_c}%
\MPIbindindexuse{MPI\_Imrecv\_c}%
\MPIbindindexuse{MPI\_Send\_init\_c}%
\MPIbindindexuse{MPI\_Bsend\_init\_c}%
\MPIbindindexuse{MPI\_Ssend\_init\_c}%
\MPIbindindexuse{MPI\_Rsend\_init\_c}%
\MPIbindindexuse{MPI\_Recv\_init\_c}%
\MPIbindindexuse{MPI\_Type\_contiguous\_c}%
\MPIbindindexuse{MPI\_Type\_vector\_c}%
\MPIbindindexuse{MPI\_Type\_create\_hvector\_c}%
\MPIbindindexuse{MPI\_Type\_indexed\_c}%
\MPIbindindexuse{MPI\_Type\_create\_hindexed\_c}%
\MPIbindindexuse{MPI\_Type\_create\_indexed\_block\_c}%
\MPIbindindexuse{MPI\_Type\_create\_hindexed\_block\_c}%
\MPIbindindexuse{MPI\_Type\_create\_struct\_c}%
\MPIbindindexuse{MPI\_Type\_create\_subarray\_c}%
\MPIbindindexuse{MPI\_Type\_create\_darray\_c}%
\MPIbindindexuse{MPI\_Type\_size\_c}%
\MPIbindindexuse{MPI\_Type\_get\_extent\_c}%
\MPIbindindexuse{MPI\_Type\_create\_resized\_c}%
\MPIbindindexuse{MPI\_Type\_get\_true\_extent\_c}%
\MPIbindindexuse{MPI\_Get\_elements\_c}%
\MPIbindindexuse{MPI\_Type\_get\_envelope\_c}%
\MPIbindindexuse{MPI\_Type\_get\_contents\_c}%
\MPIbindindexuse{MPI\_Pack\_c}%
\MPIbindindexuse{MPI\_Unpack\_c}%
\MPIbindindexuse{MPI\_Pack\_size\_c}%
\MPIbindindexuse{MPI\_Pack\_external\_c}%
\MPIbindindexuse{MPI\_Unpack\_external\_c}%
\MPIbindindexuse{MPI\_Pack\_external\_size\_c}%
\MPIbindindexuse{MPI\_Bcast\_c}%
\MPIbindindexuse{MPI\_Gather\_c}%
\MPIbindindexuse{MPI\_Gatherv\_c}%
\MPIbindindexuse{MPI\_Scatter\_c}%
\MPIbindindexuse{MPI\_Scatterv\_c}%
\MPIbindindexuse{MPI\_Allgather\_c}%
\MPIbindindexuse{MPI\_Allgatherv\_c}%
\MPIbindindexuse{MPI\_Alltoall\_c}%
\MPIbindindexuse{MPI\_Alltoallv\_c}%
\MPIbindindexuse{MPI\_Alltoallw\_c}%
\MPIbindindexuse{MPI\_Reduce\_c}%
\MPIbindindexuse{MPI\_Op\_create\_c}%
\MPIbindindexuse{MPI\_Allreduce\_c}%
\MPIbindindexuse{MPI\_Reduce\_local\_c}%
\MPIbindindexuse{MPI\_Reduce\_scatter\_block\_c}%
\MPIbindindexuse{MPI\_Reduce\_scatter\_c}%
\MPIbindindexuse{MPI\_Scan\_c}%
\MPIbindindexuse{MPI\_Exscan\_c}%
\MPIbindindexuse{MPI\_Ibcast\_c}%
\MPIbindindexuse{MPI\_Igather\_c}%
\MPIbindindexuse{MPI\_Igatherv\_c}%
\MPIbindindexuse{MPI\_Iscatter\_c}%
\MPIbindindexuse{MPI\_Iscatterv\_c}%
\MPIbindindexuse{MPI\_Iallgather\_c}%
\MPIbindindexuse{MPI\_Iallgatherv\_c}%
\MPIbindindexuse{MPI\_Ialltoall\_c}%
\MPIbindindexuse{MPI\_Ialltoallv\_c}%
\MPIbindindexuse{MPI\_Ialltoallw\_c}%
\MPIbindindexuse{MPI\_Ireduce\_c}%
\MPIbindindexuse{MPI\_Iallreduce\_c}%
\MPIbindindexuse{MPI\_Ireduce\_scatter\_block\_c}%
\MPIbindindexuse{MPI\_Ireduce\_scatter\_c}%
\MPIbindindexuse{MPI\_Iscan\_c}%
\MPIbindindexuse{MPI\_Iexscan\_c}%
\MPIbindindexuse{MPI\_Bcast\_init\_c}%
\MPIbindindexuse{MPI\_Gather\_init\_c}%
\MPIbindindexuse{MPI\_Gatherv\_init\_c}%
\MPIbindindexuse{MPI\_Scatter\_init\_c}%
\MPIbindindexuse{MPI\_Scatterv\_init\_c}%
\MPIbindindexuse{MPI\_Allgather\_init\_c}%
\MPIbindindexuse{MPI\_Allgatherv\_init\_c}%
\MPIbindindexuse{MPI\_Alltoall\_init\_c}%
\MPIbindindexuse{MPI\_Alltoallv\_init\_c}%
\MPIbindindexuse{MPI\_Alltoallw\_init\_c}%
\MPIbindindexuse{MPI\_Reduce\_init\_c}%
\MPIbindindexuse{MPI\_Allreduce\_init\_c}%
\MPIbindindexuse{MPI\_Reduce\_scatter\_block\_init\_c}%
\MPIbindindexuse{MPI\_Reduce\_scatter\_init\_c}%
\MPIbindindexuse{MPI\_Scan\_init\_c}%
\MPIbindindexuse{MPI\_Exscan\_init\_c}%
\MPIbindindexuse{MPI\_Neighbor\_allgather\_c}%
\MPIbindindexuse{MPI\_Neighbor\_allgatherv\_c}%
\MPIbindindexuse{MPI\_Neighbor\_alltoall\_c}%
\MPIbindindexuse{MPI\_Neighbor\_alltoallv\_c}%
\MPIbindindexuse{MPI\_Neighbor\_alltoallw\_c}%
\MPIbindindexuse{MPI\_Ineighbor\_allgather\_c}%
\MPIbindindexuse{MPI\_Ineighbor\_allgatherv\_c}%
\MPIbindindexuse{MPI\_Ineighbor\_alltoall\_c}%
\MPIbindindexuse{MPI\_Ineighbor\_alltoallv\_c}%
\MPIbindindexuse{MPI\_Ineighbor\_alltoallw\_c}%
\MPIbindindexuse{MPI\_Neighbor\_allgather\_init\_c}%
\MPIbindindexuse{MPI\_Neighbor\_allgatherv\_init\_c}%
\MPIbindindexuse{MPI\_Neighbor\_alltoall\_init\_c}%
\MPIbindindexuse{MPI\_Neighbor\_alltoallv\_init\_c}%
\MPIbindindexuse{MPI\_Neighbor\_alltoallw\_init\_c}%
\MPIbindindexuse{MPI\_Win\_create\_c}%
\MPIbindindexuse{MPI\_Win\_allocate\_c}%
\MPIbindindexuse{MPI\_Win\_allocate\_shared\_c}%
\MPIbindindexuse{MPI\_Win\_shared\_query\_c}%
\MPIbindindexuse{MPI\_Put\_c}%
\MPIbindindexuse{MPI\_Get\_c}%
\MPIbindindexuse{MPI\_Accumulate\_c}%
\MPIbindindexuse{MPI\_Get\_accumulate\_c}%
\MPIbindindexuse{MPI\_Rput\_c}%
\MPIbindindexuse{MPI\_Rget\_c}%
\MPIbindindexuse{MPI\_Raccumulate\_c}%
\MPIbindindexuse{MPI\_Rget\_accumulate\_c}%
\MPIbindindexuse{MPI\_File\_read\_at\_c}%
\MPIbindindexuse{MPI\_File\_read\_at\_all\_c}%
\MPIbindindexuse{MPI\_File\_write\_at\_c}%
\MPIbindindexuse{MPI\_File\_write\_at\_all\_c}%
\MPIbindindexuse{MPI\_File\_iread\_at\_c}%
\MPIbindindexuse{MPI\_File\_iread\_at\_all\_c}%
\MPIbindindexuse{MPI\_File\_iwrite\_at\_c}%
\MPIbindindexuse{MPI\_File\_iwrite\_at\_all\_c}%
\MPIbindindexuse{MPI\_File\_read\_c}%
\MPIbindindexuse{MPI\_File\_read\_all\_c}%
\MPIbindindexuse{MPI\_File\_write\_c}%
\MPIbindindexuse{MPI\_File\_write\_all\_c}%
\MPIbindindexuse{MPI\_File\_iread\_c}%
\MPIbindindexuse{MPI\_File\_iread\_all\_c}%
\MPIbindindexuse{MPI\_File\_iwrite\_c}%
\MPIbindindexuse{MPI\_File\_iwrite\_all\_c}%
\MPIbindindexuse{MPI\_File\_read\_shared\_c}%
\MPIbindindexuse{MPI\_File\_write\_shared\_c}%
\MPIbindindexuse{MPI\_File\_iread\_shared\_c}%
\MPIbindindexuse{MPI\_File\_iwrite\_shared\_c}%
\MPIbindindexuse{MPI\_File\_read\_ordered\_c}%
\MPIbindindexuse{MPI\_File\_write\_ordered\_c}%
\MPIbindindexuse{MPI\_File\_read\_at\_all\_begin\_c}%
\MPIbindindexuse{MPI\_File\_write\_at\_all\_begin\_c}%
\MPIbindindexuse{MPI\_File\_read\_all\_begin\_c}%
\MPIbindindexuse{MPI\_File\_write\_all\_begin\_c}%
\MPIbindindexuse{MPI\_File\_read\_ordered\_begin\_c}%
\MPIbindindexuse{MPI\_File\_write\_ordered\_begin\_c}%
\MPIbindindexuse{MPI\_File\_get\_type\_extent\_c}%
\MPIbindindexuse{MPI\_Register\_datarep\_c}%
  in C and through function overloading
  in the Fortran \code{mpi\_f08} module,
  (with the exception of the explicit Fortran procedures
   \MPIBCfunc{MPI\_Op\_create\_c} and \MPIBCfunc{MPI\_Register\_datarep\_c})
  and the new large count callbacks
\mpicallback{MPI\_User\_function\_c} and
\mpicallback{MPI\_Datarep\_conversion\_function\_c}
together with the predefined function
\mpiconstfunc{MPI\_CONVERSION\_FN\_NULL\_C}
  were introduced to accomodate large buffers and/or datatypes.

  Clarifications were added to the behavior of
  \INOUT/\OUT\ parameters that cannot represent the value
  to be returned for the \mpifunc{MPI\_BUFFER\_DETACH} and
  \mpifunc{MPI\_FILE\_GET\_TYPE\_EXTENT} functions.

  A new error class \error{MPI\_ERR\_VALUE\_TOO\_LARGE} was
  introduced.

% 04.--- MPI-4.0 Issues 1, 3 p24
\item
    Sections~\ref{sec:terms-errorhandling},~\ref{sec:errorhandler},~\ref{sec:ei-error}, and~\ref{sec:inquiry-startup} on
pages~\pageref{sec:terms-errorhandling},~\pageref{sec:errorhandler},~\pageref{sec:ei-error}, and~\pageref{sec:inquiry-startup}.
\newline
\MPI/ calls that are not related to any objects are considered to be attached to
the communicator \mpiconst{MPI\_COMM\_SELF} instead of \mpiconst{MPI\_COMM\_WORLD}.
The definition of
\mpiconst{MPI\_ERRORS\_ARE\_FATAL} was clarified to cover all connected processes,
and a new error handler, \mpiconst{MPI\_ERRORS\_ABORT}, was created to limit the
scope of aborting.

% 05.--- MPI-4.0-issue #122 p58
\item
\mpitermindex{cancel}
Section~\ref{sec:pt2pt-nonblock} on page~\pageref{sec:pt2pt-nonblock}.
\newline
The introduction of \MPI/ nonblocking communication was changed to
describe correctness and performance reasons for the use of nonblocking communication.

% 06.--- MPI-4.0-issue #125 p60
\item
  Section~\ref{subsec:pt2pt-commstart} on page~\pageref{subsec:pt2pt-commstart}.
  \newline
  Addition of
  \mpifunc{MPI\_ISENDRECV} and
  \mpifunc{MPI\_ISENDRECV\_REPLACE}.

% 07.--- MPI-4.0-issue #25 p67
\item
Sections~\ref{subsec:pt2pt-commend}, \ref{sec:pt2pt-persistent}, \ref{sec:coll-persistent}, \ref{sec:sparsecoll-pnbc}, and \ref{sec:topol-applic-example} on pages~\pageref{subsec:pt2pt-commend}, \pageref{sec:pt2pt-persistent}, \pageref{sec:coll-persistent}, \pageref{sec:sparsecoll-pnbc}, and \pageref{sec:topol-applic-example}.
\newline
Persistent collective communication 
\mpifuncindex{MPI\_ALLGATHER\_INIT}%
\mpifuncindex{MPI\_ALLGATHERV\_INIT}%
\mpifuncindex{MPI\_ALLREDUCE\_INIT}%
\mpifuncindex{MPI\_ALLTOALL\_INIT}%
\mpifuncindex{MPI\_ALLTOALLV\_INIT}%
\mpifuncindex{MPI\_ALLTOALLW\_INIT}%
\mpifuncindex{MPI\_BARRIER\_INIT}%
\mpifuncindex{MPI\_BCAST\_INIT}%
\mpifuncindex{MPI\_EXSCAN\_INIT}%
\mpifuncindex{MPI\_GATHER\_INIT}%
\mpifuncindex{MPI\_GATHERV\_INIT}%
\mpifuncindex{MPI\_REDUCE\_INIT}%
\mpifuncindex{MPI\_REDUCE\_SCATTER\_BLOCK\_INIT}%
\mpifuncindex{MPI\_REDUCE\_SCATTER\_INIT}%
\mpifuncindex{MPI\_SCAN\_INIT}%
\mpifuncindex{MPI\_SCATTER\_INIT}%
\mpifuncindex{MPI\_SCATTERV\_INIT}%
\mpiskipfunc{MPI\_\{ALLGATHER$|$\ldots\}\_INIT}
including persistent collective neighborhood communication
\mpifuncindex{MPI\_NEIGHBOR\_ALLGATHER\_INIT}%
\mpifuncindex{MPI\_NEIGHBOR\_ALLGATHERV\_INIT}%
\mpifuncindex{MPI\_NEIGHBOR\_ALLTOALL\_INIT}%
\mpifuncindex{MPI\_NEIGHBOR\_ALLTOALLV\_INIT}%
\mpifuncindex{MPI\_NEIGHBOR\_ALLTOALLW\_INIT}%
\vlgb\mpiskipfunc{MPI\_NEIGHBOR\_\{ALLGATHER$|$\ldots\}\_INIT}
was added to the standard.

% 08.--- MPI-4.0-issue #26 p88
\item
\mpitermindex{cancel}
Sections~\ref{sec:cancel} and~\ref{sec:deprecated:since40} on pages~\pageref{sec:cancel} and~\pageref{sec:deprecated:since40}.
\newline
Cancelling a send request by calling \mpifunc{MPI\_CANCEL} has been deprecated and may be removed in a future version of the \MPI/ specification.

% 09.--- MPI-4.0-issue #136 p99
\item
  Chapter~\ref{chap:part} on page~\pageref{chap:part}.
  \newline
  A new chapter on partitioned communication 
  with the new \MPI/ procedures 
\mpifuncindex{MPI\_PARRIVED}%
\mpifuncindex{MPI\_PREADY}%
\mpifuncindex{MPI\_PREADY\_LIST}%
\mpifuncindex{MPI\_PREADY\_RANGE}%
\mpifuncindex{MPI\_PRECV\_INIT}%
\mpifuncindex{MPI\_PSEND\_INIT}%
  \mpiskipfunc{MPI\_\{PARRIVED$|$PREADY\{\ldots\}\}}
  and \mpiskipfunc{MPI\_\{PRECV$|$PSEND\}\_INIT}
  was added.

% 10.--- MPI-4.0-issue #156 p321
\item
Section~\ref{subsec:context-intracomconst} on page~\pageref{subsec:context-intracomconst}.
\newline
\mpiconst{MPI\_COMM\_TYPE\_HW\_UNGUIDED} was added as a new possible value for the \mpiarg{split\_type}
  parameter of the \mpifunc{MPI\_COMM\_SPLIT\_TYPE} function.

% 11.--- MPI-4.0-issue #132 p321
\item
  Section~\ref{subsec:context-intracomconst} on page~\pageref{subsec:context-intracomconst}.
  \newline
  \mpiconst{MPI\_COMM\_TYPE\_HW\_GUIDED} was added as a new possible value for the \mpiarg{split\_type}
  parameter of the \mpifunc{MPI\_COMM\_SPLIT\_TYPE} function, as well as a new info key
  \infokey{mpi\_hw\_resource\_type}. A specific value associated with this new info key is also defined:
  \infoval{mpi\_shared\_memory}.

% 12.--- MPI-4.0-issue #52 p321
\item
Section~\ref{subsec:context-intracomconst} on page~\pageref{subsec:context-intracomconst}.
\newline
The functions \mpifunc{MPI\_COMM\_DUP} and \mpifunc{MPI\_COMM\_IDUP} were
updated to no longer propagate info hints.
\newline This change may affect backward compatibility.

% 13.--- MPI-4.0-issue #53 p321
\item
Section~\ref{subsec:context-intracomconst} on page~\pageref{subsec:context-intracomconst}.
\newline
The
\mpifunc{MPI\_COMM\_IDUP\_WITH\_INFO} function was added.

% 14.--- MPI-4.0-issue #11 p339.1
\item
Sections~\ref{sec:comm-info}, \ref{subsec:window-info}, and~\ref{sec:io-info} on pages
\pageref{sec:comm-info}, \pageref{subsec:window-info}, and \pageref{sec:io-info}.
\newline
The definition of info hints was updated to allow applications to provide
assertions regarding their usage of \MPI/ objects and operations.

% 15.--- MPI-4.0-issue #11 p339.2
\item
Section~\ref{sec:comm-info} on page~\pageref{sec:comm-info}.
\newline
The new info hints \infokey{mpi\_assert\_no\_any\_tag},
\infokey{mpi\_assert\_no\_any\_source}, \infokey{mpi\_assert\_exact\_length},
and \infokey{mpi\_assert\_allow\_overtaking} were added for use with
communicators.

% 16.--- MPI-4.0-issue #63 p339
\item
Sections~\ref{sec:comm-info}, \ref{subsec:window-info}, and~\ref{sec:io-info}
on pages~\pageref{sec:comm-info}, \pageref{subsec:window-info},
and~\pageref{sec:io-info}.
\newline
The semantics of the
\mpifunc{MPI\_COMM\_SET\_INFO}, \mpifunc{MPI\_COMM\_GET\_INFO},
\mpifunc{MPI\_WIN\_SET\_INFO}, \mpifunc{MPI\_WIN\_GET\_INFO},
\mpifunc{MPI\_FILE\_SET\_INFO}, and \mpifunc{MPI\_FILE\_GET\_INFO}
were clarified.

% 17.--- MPI-4.0-issue #72 p387
\item
Section~\ref{subsec:topol-construct} on page~\pageref{subsec:topol-construct}.
\newline
\mpifunc{MPI\_DIMS\_CREATE} is now guaranteed to return \error{MPI\_SUCCESS} if the number of dimensions passed to the routine is set to 0 and the number of nodes is set to 1.

% 18.--- MPI-4.0-issue #121 p447
\mpitermindex{memory!alignment}\mpitermindex{alignment}
\item Sections~\ref{sec:misc-memalloc},~\ref{sec:winalloc}, and~\ref{sec:winallocshared} on
pages~\pageref{sec:misc-memalloc},~\pageref{sec:winalloc}, and~\pageref{sec:winallocshared}.
\newline
{\raggedright
Introduced alignment requirements for memory allocated through
\mpifunc{MPI\_ALLOC\_MEM}, \mpifunc{MPI\_WIN\_ALLOCATE}, and
\mpifunc{MPI\_WIN\_ALLOCATE\_SHARED} and added a new \mpiarg{info}
key \infokey{mpi\_minimum\_memory\_alignment} to specify a desired alternative minimum alignment.

}
% 19.--- MPI-4.0-issue #28 p450
\item
Sections~\ref{sec:errorhandler} and~\ref{sec:ei-error-classes} on
pages~\pageref{sec:errorhandler} and~\pageref{sec:ei-error-classes}.
\newline
Clarified definition of errors to say that \MPI/ should continue whenever
possible and allow the user to recover from errors.

% 20.--- MPI-4.0-issue #119 p461
\item
Section~\ref{sec:ei-error-classes} on page~\pageref{sec:ei-error-classes}.
\newline
Added text to clarify what is implied about the status of \MPI/ and user visible buffers when \MPI/
functions return \error{MPI\_SUCCESS} or other error codes.

% 21.---  MPI-4.0-issue #134 p463
\item
  Section~\ref{sec:ei-error-classes} on page~\pageref{table:inquiry:errclasses:part:i}.
  \newline
  The error class \error{MPI\_ERR\_PROC\_ABORTED} has been added.

% 22.--- MPI-4.0-issue #129 p471
%% label sec:info has become subsec:info. If the original number in the
%% MPI-3 document, these should use the exact number, not a labeled referece.
\item
Section~\ref{subsec:info} on page~\pageref{subsec:info}.
\newline
Added a new function \mpifunc{MPI\_INFO\_GET\_STRING} that takes a
buffer length argument for returning info value strings. This function
returns the required buffer length for the requested string and guarantees
null termination for C strings where buffer size is greater than 0.

% 23.--- MPI-4.0-issue #131 p471
\item
Section~\ref{subsec:info} on page~\pageref{subsec:info} and
Section~\ref{sec:deprecated:since40} on page~\pageref{sec:deprecated:since40}.
\newline
\mpifunc{MPI\_INFO\_GET} and \mpifunc{MPI\_INFO\_GET\_VALUELEN} were deprecated.
  
% 24.--- MPI-4.0-issue #103 p479
\item
Chapter~\ref{sec:dynamic-2},
\ref{subsec:pt2pt-envelope},
\ref{sec:predef-comms},
\ref{subsec:context-grpconst},
\ref{subsec:context-intracomconst},
\ref{subsec:context-intercomm},
\ref{subsec:inquiry-version},
\ref{subsec:inquiry-inquiry},
\ref{sec:errorhandler},
\ref{subsec:inquiry-errhdlr-sessions},
\ref{sec:ei-error},
\ref{sec:ei-threads},
\ref{sec:io-filecntl-open},
\ref{sec:io-filecntl-get-params},
\ref{sec:io-errhandlers},
\ref{sec:mpit:init},
\ref{sec:misc-handleconvert},
\ref{subsec:mpiopaqueobjects},
and
Annex~\ref{sec:lang}
on pages~\pageref{sec:dynamic-2},
\pageref{subsec:pt2pt-envelope},
\pageref{sec:predef-comms},
\pageref{subsec:context-grpconst},
\pageref{subsec:context-intracomconst},
\pageref{subsec:context-intercomm},
\pageref{subsec:inquiry-version},
\pageref{subsec:inquiry-inquiry},
\pageref{sec:errorhandler},
\pageref{subsec:inquiry-errhdlr-sessions},
\pageref{sec:ei-error},
\pageref{sec:ei-threads},
\pageref{sec:io-filecntl-open},
\pageref{sec:io-filecntl-get-params},
\pageref{sec:io-errhandlers},
\pageref{sec:mpit:init},
\pageref{sec:misc-handleconvert},
\pageref{subsec:mpiopaqueobjects},
and
\pageref{sec:lang}.
\newline
{\raggedright
The \spm{} was added to the standard.
New \MPI/ procedures are 
\mpifuncindex{MPI\_SESSION\_INIT}%
\mpifuncindex{MPI\_SESSION\_FINALIZE}%
\mpiskipfunc{MPI\_SESSION\_\{INIT$|$FINALIZE\}},
\mpifuncindex{MPI\_SESSION\_GET\_INFO}%
\mpifuncindex{MPI\_SESSION\_GET\_NTH\_PSET}%
\mpifuncindex{MPI\_SESSION\_GET\_NUM\_PSETS}%
\mpifuncindex{MPI\_SESSION\_GET\_PSET\_INFO}%
\mpiskipfunc{MPI\_SESSION\_GET\_\{\ldots\}},
\mpifuncindex{MPI\_SESSION\_CREATE\_ERRHANDLER}%
\mpifuncindex{MPI\_SESSION\_SET\_ERRHANDLER}%
\mpifuncindex{MPI\_SESSION\_GET\_ERRHANDLER}%
\mpifuncindex{MPI\_SESSION\_CALL\_ERRHANDLER}%
\mpiskipfunc{MPI\_SESSION\_\{\ldots\}\_ERRHANDLER},
\mpifunc{MPI\_GROUP\_FROM\_SESSION\_PSET}, 
\mpifunc{MPI\_COMM\_CREATE\_FROM\_GROUP}, 
\vlgb\mpifunc{MPI\_INTERCOMM\_CREATE\_FROM\_GROUPS},
\mpifuncindex{MPI\_SESSION\_C2F}%
\mpifuncindex{MPI\_SESSION\_F2C}%
and new conversion functions are \mpiskipfunc{MPI\_SESSION\_\{C2F$|$F2C\}}.
New declarations are \mpictype{MPI\_Session} in C 
and \mpiftype{TYPE(MPI\_Session)}
together with the related overloaded operators \code{.EQ.}, \code{.NE.}, \code{==} and \code{/=}
in the Fortran \code{mpi\_f08} and \code{mpi} modules,
and the callback function prototype \mpicallback{MPI\_Session\_errhandler\_function}\mpicallbackindex{SESSION\_ERRHANDLER\_FUNCTION}.
New constants are \mpiconst{MPI\_SESSION\_NULL}, \error{MPI\_ERR\_SESSION},
\mpiconst{MPI\_MAX\_PSET\_NAME\_LEN}, \mpiconst{MPI\_MAX\_STRINGTAG\_LEN}, \mpiconst{MPI\_T\_BIND\_MPI\_SESSION}
and the predefined info key \infokey{mpi\_size}.

}
% 25.--- MPI-4.0-issue #143 p480
\item
Section~\ref{sec:inquiry-startup} on page~\pageref{sec:inquiry-startup}.
\newline
A new function \mpifunc{MPI\_INFO\_CREATE\_ENV} was added.

% 26.--- MPI-4.0-issue #102 p480
\item
    Sections~\ref{sec:inquiry-startup} and~\ref{subsec:disconnect}
    on pages~\pageref{sec:inquiry-startup} and~\pageref{subsec:disconnect}.
    \newline
    Clarified the semantic of failure and error reporting before (and during) \mpifunc{MPI\_INIT} and after \mpifunc{MPI\_FINALIZE}.

% 27.--- MPI-4.0-issue #102 p521
\item
    Section~\ref{subsec:spawnkeys} on page~\pageref{subsec:spawnkeys}.
    \newline
    Added the \infokey{mpi\_initial\_errhandler} reserved info key with the
    reserved values \infoval{mpi\_errors\_abort},
    \infoval{mpi\_errors\_are\_fatal}, and \infoval{mpi\_errors\_return} to the
    launch keys in \mpifunc{MPI\_COMM\_SPAWN}, \mpifunc{MPI\_COMM\_SPAWN\_MULTIPLE}, and \mpilaunch{mpiexec}.

% 28.--- MPI-4.0-issue #405 p592
\item
  Section~\ref{sec:1sided-lock} on page~\pageref{sec:1sided-lock}.
  \newline
  \RMA/ passive target synchronization using locks can now be used portably in memory
  allocated via \mpifunc{MPI\_WIN\_ALLOCATE\_SHARED}.

% 29.--- MPI-4.0-issue #91 p628
% No label mpi:chap:ei:status-set-cancelled
\item
  Section~\ref{sec:ei-status} on page~\pageref{sec:ei-status}.
  \newline
  The \code{mpi\_f08} binding incorrectly had the dummy parameter
  \code{flag} in the \MPI/ F08 binding for
  \mpifunc{MPI\_STATUS\_SET\_CANCELLED} marked as
  \code{INTENT(OUT)}.  It has been fixed to be \code{INTENT(IN)}.

% 30.--- MPI_4.0-issue #79 p745
\item
Sections~\ref{sec:mpit} and~\ref{sec:mpit:events} on pages~\pageref{sec:mpit} and~\pageref{sec:mpit:events}.
\newline
{\raggedright
A callback-driven event interface
\mpifuncindex{MPI\_T\_CATEGORY\_GET\_EVENTS}%
\mpifuncindex{MPI\_T\_CATEGORY\_GET\_NUM\_EVENTS}%
\mpifuncindex{MPI\_T\_EVENT\_CALLBACK\_GET\_INFO}%
\mpifuncindex{MPI\_T\_EVENT\_CALLBACK\_SET\_INFO}%
\mpifuncindex{MPI\_T\_EVENT\_COPY}%
\mpifuncindex{MPI\_T\_EVENT\_GET\_INDEX}%
\mpifuncindex{MPI\_T\_EVENT\_GET\_INFO}%
\mpifuncindex{MPI\_T\_EVENT\_GET\_NUM}%
\mpifuncindex{MPI\_T\_EVENT\_GET\_SOURCE}%
\mpifuncindex{MPI\_T\_EVENT\_GET\_TIMESTAMP}%
\mpifuncindex{MPI\_T\_EVENT\_HANDLE\_ALLOC}%
\mpifuncindex{MPI\_T\_EVENT\_HANDLE\_FREE}%
\mpifuncindex{MPI\_T\_EVENT\_HANDLE\_GET\_INFO}%
\mpifuncindex{MPI\_T\_EVENT\_HANDLE\_SET\_INFO}%
\mpifuncindex{MPI\_T\_EVENT\_READ}%
\mpifuncindex{MPI\_T\_EVENT\_REGISTER\_CALLBACK}%
\mpifuncindex{MPI\_T\_EVENT\_SET\_DROPPED\_HANDLER}%
\mpifuncindex{MPI\_T\_SOURCE\_GET\_INFO}%
\mpifuncindex{MPI\_T\_SOURCE\_GET\_NUM}%
\mpifuncindex{MPI\_T\_SOURCE\_GET\_TIMESTAMP}%
with the \mpiskipfunc{MPI\_T\_\{SOURCE$|$EVENT\}\_\{\ldots\}}
and \mpiskipfunc{MPI\_T\_CATEGORY\_\{GET$|$GET\_NUM\}\_EVENTS} routines,
the declaration types \mpictype{MPI\_T\_cb\_safety},
\mpictypeskip{MPI\_T\_event\_\{instance$|$registration\}}\cdeclindex{MPI\_T\_event\_instance}\cdeclindex{MPI\_T\_event\_registration},
\mpictype{MPI\_T\_source\_order},
and the callback function prototypes \mpicallbackskip{MPI\_T\_event\_\{cb$|$dropped\_cb$|$free\_cb\}\_function},
\mpicallbackindex{MPI\_T\_event\_cb\_function}%
\mpicallbackindex{MPI\_T\_event\_dropped\_cb\_function}%
\mpicallbackindex{MPI\_T\_event\_free\_cb\_function}%
were added to the \MPI/ tool information interface.

}
% 31.--- MPI-4.0-issue #392 p772
\item
  Section~\ref{sec:mpit:category_query_functions} on page~\pageref{sec:mpit:category_changed}.
  \newline
  The argument \mpiarg{stamp} (previously described as a virtual time stamp) from \mpifunc{MPI\_T\_CATEGORY\_CHANGED} was renamed to \mpiarg{update\_number} and its intended implementation and use was clarified.

% 32.--- MPI-4.0-issue #60 p766
\item
Section~\ref{sec:mpit:retcodes}, Table~\ref{table:tools:mpit:err}, and Section~\ref{sec:deprecated:since40}
on pages~\pageref{sec:mpit:retcodes}, \pageref{table:tools:mpit:err}, and~\pageref{sec:deprecated:since40}.
\newline
\error{MPI\_T\_ERR\_INVALID\_ITEM} is deprecated. \MPI/
routines should return
\newline
\error{MPI\_T\_ERR\_INVALID\_INDEX} instead of
\error{MPI\_T\_ERR\_INVALID\_ITEM}.

% 33.--- MPI-4.0-issue #51 p773
\item
  Section~\ref{sec:deprecated:since40} on page~\pageref{sec:deprecated-fortran-sizeof}.
  \newline
  \mpifunc{MPI\_SIZEOF} was deprecated.

% 34.--- MPI-4.0-issue #405 p786
\item
  Section~\ref{sec:f90:linker-names} on page~\pageref{sec:f90:linker-names}.
  \newline
  An exception was added for the specific Fortran names in the case of TS 29113 interface specifications 
  in \code{mpif.h} for \mpifunc{MPI\_NEIGHBOR\_ALLTOALLW\_INIT}, 
  \mpifunc{MPI\_NEIGHBOR\_ALLTOALLV\_INIT}, and \mpifunc{MPI\_NEIGHBOR\_ALLGATHERV\_INIT}.

\end{enumerate}

\section{Changes from Version 3.0 to Version 3.1}
\label{subsec:30to31}
% The following tickets do not have a change-log entry:
% 
% 1) The following tickets are passed MPI-3.0-errata tickets:
% 
%    Tickets that need NOT a change-log and do not have one:
%      #347 #348 #354 #359 #367 #385 #387 #390 #393 #403 #405
%      #406 #413 #420 #431 #433 #440 #441 #443 #444 #445 #446
%      #447 #449 #450 #451 #452 #453
% 
%    Tickets that MAY need a change-log and do NOT yet have one:
%      #383 #394
% 
%    Tickets that need a change-log and do NOT yet have one:
%      #345 #415
%
%    Tickets that need a change-log and have one:
%      #350 #355 #362 #386 #388 #389 #391 #419 #424
%      (#421 belongs to #349, see below)
% 
% 2) The following tickets are passed MPI-3.1 tickets:
% 
%    Tickets that need NOT a change-log and do not have one:
%      #368 #455
% 
%    Tickets that need a change-log and do NOT yet have one:
%      #357 #400
%
%    Tickets that need a change-log and have one:
%      #273 #349+#402+#404+#421 #369 #377 #378

\subsection{Fixes to Errata in Previous Versions of \texorpdfstring{\MPI/}{MPI}} 
\label{sec:changes-30-31-errata}
 
\begin{enumerate}

% 01.--- MPI-3.0-erratum Ticket 388
\item 
Chapters~\ref{chap:pt2pt}--\ref{chap:binding-2},
Annex~\ref{sec:lang:fortran-mpi-f08-module} on page~\pageref{sec:lang:fortran-mpi-f08-module}, 
and Example~\ref{coll-exX-fortran} on page~\pageref{coll-exX-fortran}, and 
\MPIIIIDOTO/ Chapters 3--17, Annex A.3 on page 707, and Example 5.21 on page 187.
\newline 
Within the \code{mpi\_f08} Fortran support method, \ftype{BIND(C)} was removed from 
all \flushline
\code{SUBROUTINE}, \code{FUNCTION}, and \code{ABSTRACT INTERFACE} definitions.

% 02.--- MPI-3.0-erratum Ticket 415
\item 
Section~\ref{subsec:pt2pt-status} on page~\pageref{subsec:pt2pt-status}, and
\MPIIIIDOTO/ Section~3.2.5 on page~30.
\newline 
The three public fields \mpiconst{MPI\_SOURCE}, \mpiconst{MPI\_TAG}, and \mpiconst{MPI\_ERROR} 
of the Fortran derived type \mpiftype{TYPE(MPI\_Status)}\cdeclindex{MPI\_Status}
must be of type \ftype{INTEGER}.

% 03.--- MPI-3.0-erratum Ticket 424
\item 
Section~\ref{sec:matching-probe} on page~\pageref{sec:matching-probe}, and
\MPIIIIDOTO/ Section~3.8.2 on page~67.
\newline 
The flag arguments of the Fortran interfaces of \mpifunc{MPI\_IMPROBE} 
were originally incorrectly defined as \ftype{INTEGER} (instead as \ftype{LOGICAL}). 

% 04.--- MPI-3.0-erratum Ticket 345
\item 
Section~\ref{subsec:context-intracomconst} on page~\pageref{subsec:context-intracomconst}, and
\MPIIIIDOTO/ Section~6.4.2 on page 237.
\newline 
In the \code{mpi\_f08} binding of \mpifunc{MPI\_COMM\_IDUP},
the output argument \mpiarg{newcomm} is declared as \ftype{ASYNCHRONOUS}.

% 05.--- MPI-3.0-erratum Ticket 345
\item 
Section~\ref{sec:comm-info} on page~\pageref{sec:comm-info}, and
\MPIIIIDOTO/ Section~6.4.4 on page 248.
\newline 
In the \code{mpi\_f08} binding of \mpifunc{MPI\_COMM\_SET\_INFO},
the \code{INTENT} of \mpiarg{comm} is \code {IN}, and
the optional output argument \mpiarg{ierror} was missing.

% 06.--- MPI-3.0-erratum Ticket 419
\item 
Section~\ref{sec:sparsecoll} on page~\pageref{sec:sparsecoll}, and
\MPIIIIDOTO/ Sections 7.6, on pages 314.
\newline 
In the case of virtual general graph topolgies (created with \mpifunc{MPI\_CART\_CREATE}),
the use of neighborhood collective communication is restricted to adjacency matrices
with the number of edges between any two processes is defined to be the same for
both processes (i.e., with a symmetric adjacency matrix). 

% 07.--- MPI-3.0-erratum Ticket 345
\item 
Section~\ref{subsec:inquiry-version} on page~\pageref{subsec:inquiry-version}, and
\MPIIIIDOTO/ Section~8.1.1 on page 335.
\newline 
In the \code{mpi\_f08} binding of \mpifunc{MPI\_GET\_LIBRARY\_VERSION},
a typo in the \mpiarg{resultlen} argument was corrected.

% 08.--- MPI-3.0-erratum Ticket 388
\item 
{\raggedright
Sections~\ref{sec:misc-memalloc} (\mpifunc{MPI\_ALLOC\_MEM} and \mpifunc{MPI\_ALLOC\_MEM\_CPTR}),\\
\ref{sec:winalloc} (\mpifunc{MPI\_WIN\_ALLOCATE} and \mpifunc{MPI\_WIN\_ALLOCATE\_CPTR}),\\
\ref{sec:winallocshared} (\mpifunc{MPI\_WIN\_ALLOCATE\_SHARED} and \mpifunc{MPI\_WIN\_ALLOCATE\_SHARED\_CPTR}),\\
\ref{sec:winallocshared} (\mpifunc{MPI\_WIN\_SHARED\_QUERY} and \mpifunc{MPI\_WIN\_SHARED\_QUERY\_CPTR}),\\
\ref{sec:prof:require} and~\ref{sec:prof:complications} (Profiling interface), 
and corresponding sections in \MPIIIIDOTO/.
\newline 
The linker name concept was substituted by defining specific procedure names.

}
% 09.--- MPI-3.0-erratum Ticket 362
\item 
Section~\ref{chap:one-side-2:win_create} on page~\pageref{chap:one-side-2:win_create}, and
\MPIIIIDOTO/ Section~11.2.2 on page 407.
\newline 
The \infokey{same\_size} info key can be used with all window flavors,
and requires that all processes in the process group of the communicator
have provided this info key with the same value.


% 10.--- MPI-3.0-erratum Ticket 350
\item 
Section~\ref{sec:1sided-accumulate} on page~\pageref{sec:1sided-accumulate}, and
\MPIIIIDOTO/ Section~11.3.4 on page 424.
\newline 
Origin buffer arguments to \mpifunc{MPI\_GET\_ACCUMULATE} are ignored when the \mpiconst{MPI\_NO\_OP} operation is used. 

% 11.--- MPI-3.0-erratum Ticket 355
\item 
Section~\ref{sec:1sided-accumulate} on page~\pageref{sec:1sided-accumulate}, and
\MPIIIIDOTO/ Section~11.3.4 on page 424.
\newline 
Clarify the roles of origin, result, and target communication parameters in \mpifunc{MPI\_GET\_ACCUMULATE}. 

% 12.--- MPI-3.0-erratum Ticket 383
\item 
Section~\ref{sec:mpit} on page~\pageref{sec:mpit}, and
\MPIIIIDOTO/ Section 14.3 on page~561
\newline 
New paragraph and advice to users clarifying intent of variable names in the tools information interface.

% 13.--- MPI-3.0-erratum Ticket 383
\item 
Section~\ref{sec:mpit:strings} on page~\pageref{sec:mpit:strings}, and
\MPIIIIDOTO/ Section~14.3.3 on page~563.
\newline 
New paragraph clarifying variable name equivalence in the tools information interface.

% 14.--- MPI-3.0-erratum Ticket 383
\item 
Sections~\ref{sec:mpit:cvar}, \ref{sec:mpit:pvar}, and~\ref{sec:mpit:cat} on pages~\pageref{sec:mpit:cvar}, \pageref{sec:mpit:pvar}, and \pageref{sec:mpit:cat}, and
\newline 
\MPIIIIDOTO/ Sections~14.3.6, 14.3.7, and 14.3.8 on pages 567, 573, and 584.
\newline 
In functions \mpifunc{MPI\_T\_CVAR\_GET\_INFO}, \mpifunc{MPI\_T\_PVAR\_GET\_INFO}, and \mpifunc{MPI\_T\_CATEGORY\_GET\_INFO}, clarification of parameters that must be identical for equivalent control variable / performance variable / category names across connected processes.

% 15.--- MPI-3.0-erratum Ticket 391
\item 
Section~\ref{sec:mpit:pvar} on page~\pageref{sec:mpit:pvar}, and
\MPIIIIDOTO/ Section~14.3.7 on page~573.
\newline 
Clarify return code 
\mpifuncindex{MPI\_T\_PVAR\_START}%
\mpifuncindex{MPI\_T\_PVAR\_STOP}%
\mpifuncindex{MPI\_T\_PVAR\_RESET}%
of \mpiskipfunc{MPI\_T\_PVAR\_\{START,STOP,RESET\}} routines. 

% 16.--- MPI-3.0-erratum Ticket 386
\item 
Section~\ref{sec:mpit:pvar} on page~\pageref{sec:mpit:pvar}, and
\MPIIIIDOTO/ Section~14.3.7 on page~579, line~7.
\newline 
Clarify the return code when bad handle is passed to 
\mpifuncindex{MPI\_T\_PVAR\_HANDLE\_FREE}%
\mpifuncindex{MPI\_T\_PVAR\_START}%
\mpifuncindex{MPI\_T\_PVAR\_STOP}%
\mpifuncindex{MPI\_T\_PVAR\_READ}%
\mpifuncindex{MPI\_T\_PVAR\_WRITE}%
\mpifuncindex{MPI\_T\_PVAR\_RESET}%
\mpifuncindex{MPI\_T\_PVAR\_READRESET}%
an \mpiskipfunc{MPI\_T\_PVAR\_*} routine. 

% 17.--- MPI-3.0-erratum Ticket 388
\item 
Section~\ref{f90:basic} on page~\pageref{f90:basic}, and
\MPIIIIDOTO/ Section~17.1.4 on page~603.
\newline 
The advice to implementors at the end of the section was rewritten and moved into the following section.

% 18.--- MPI-3.0-erratum Ticket 388
\item 
Section~\ref{sec:f90:linker-names} on page~\pageref{sec:f90:linker-names}, and
\MPIIIIDOTO/ Section~17.1.5 on page~605.
\newline 
The section was fully rewritten. The linker name concept was substituted by defining specific procedure names.

% 19.--- MPI-3.0-erratum Ticket 388
\item 
Section~\ref{sec:f90:different-fortran-versions} on page~\pageref{sec:f90:different-fortran-versions}, and
\MPIIIIDOTO/ Section~17.1.6 on page~611.
\newline 
The requirements on \ftype{BIND(C)} procedure interfaces were removed. 

% 20.--- MPI-3.0-erratum Ticket 389
\item 
Annexes~\ref{sec:lang:c-mpi-interfaces}, \ref{sec:lang:fortran-mpi-f08-module}, and~\ref{sec:lang:fortran-mpifh-and-mpi-module}
on pages~\pageref{sec:lang:c-mpi-interfaces}, \pageref{sec:lang:fortran-mpi-f08-module}, and~\pageref{sec:lang:fortran-mpifh-and-mpi-module}, and
\newline 
\MPIIIIDOTO/ Annexes~A.2, A.3, and A.4 on pages~685, 707, and 756.
\newline 
The predefined callback \mpiconstfunc{MPI\_CONVERSION\_FN\_NULL} was added to all three annexes.

% 21.--- MPI-3.0-erratum Ticket 345
\item 
{\raggedright
Annex~\ref{app:subsec:f08:context} on page~\pageref{app:subsec:f08:context}, and
\MPIIIIDOTO/ Annex~A.3.4 on page~724.
\newline 
In the \code{mpi\_f08} binding 
\mpiconstfuncindex{MPI\_COMM\_DUP\_FN}%
\mpiconstfuncindex{MPI\_COMM\_NULL\_COPY\_FN}%
\mpiconstfuncindex{MPI\_COMM\_NULL\_DELETE\_FN}%
\mpiconstfuncindex{MPI\_TYPE\_DUP\_FN}%
\mpiconstfuncindex{MPI\_TYPE\_NULL\_COPY\_FN}%
\mpiconstfuncindex{MPI\_TYPE\_NULL\_DELETE\_FN}%
\mpiconstfuncindex{MPI\_WIN\_DUP\_FN}%
\mpiconstfuncindex{MPI\_WIN\_NULL\_COPY\_FN}%
\mpiconstfuncindex{MPI\_WIN\_NULL\_DELETE\_FN}%
of\\
\mpiskipfunc{MPI\_\{COMM$|$TYPE$|$WIN\}\_\{DUP$|$NULL\_COPY$|$NULL\_DELETE\}\_FN},
all \code{INTENT(...)} information was removed.

}
\end{enumerate}

\subsection{Changes in \texorpdfstring{\MPIIIIDOTI/}{MPI-3.1}}
\label{sec:changes-30-31}
 
\begin{enumerate}

% 01.--- MPI-3.0 Tickets 349, 402, 404, 421  
\item 
Sections~\ref{sec:macros} and~\ref{subsec:pt2pt-addfunc} on pages~\pageref{sec:macros} and~\pageref{subsec:pt2pt-addfunc}.
\newline 
The use of the intrinsic operators ``\code{+}'' and ``\code{-}'' for absolute addresses is substituted by \mpifunc{MPI\_AINT\_ADD} and \mpifunc{MPI\_AINT\_DIFF}. In C, they can be implemented as macros. 

% 02.--- MPI-3.1 Ticket 357
\item 
Sections~\ref{subsec:inquiry-version}, \ref{sec:inquiry-startup}, and~\ref{sec:ei-threads}
on pages~\pageref{subsec:inquiry-version}, \pageref{sec:inquiry-startup}, and~\pageref{sec:ei-threads}.
\newline 
The routines \mpifunc{MPI\_INITIALIZED}, \mpifunc{MPI\_FINALIZED},
\mpifunc{MPI\_QUERY\_THREAD}, \mpifunc{MPI\_IS\_THREAD\_MAIN}, 
\mpifunc{MPI\_GET\_VERSION}, and \mpifunc{MPI\_GET\_LIBRARY\_VERSION}
are callable from threads without restriction (in the sense of 
\mpiconst{MPI\_THREAD\_MULTIPLE}), irrespective of the actual level of thread support 
provided, in the case where the implementation supports threads.  

% 03.--- MPI-3.1 Ticket 369
\item 
Section~\ref{chap:one-side-2:win_create} on page~\pageref{chap:one-side-2:win_create}.
\newline 
The \infokey{same\_disp\_unit} info key was added for use in \RMA/ window creation routines. 

% 04.--- MPI-3.1 Ticket 273
\item 
Sections~\ref{sec:io-explicit} and~\ref{sec:io-indiv-ptr} 
on pages~\pageref{sec:io-explicit} and~\pageref{sec:io-indiv-ptr}.
\newline 
%% WDG - Corrected to refer to the language-neutral names, as required
%% by the standard and the mpifunc macro
Added \mpifunc{MPI\_FILE\_IREAD\_AT\_ALL}, \mpifunc{MPI\_FILE\_IWRITE\_AT\_ALL},
\mpifunc{MPI\_FILE\_IREAD\_ALL}, and \mpifunc{MPI\_FILE\_IWRITE\_ALL}

% 05.--- MPI-3.1 Ticket 378
\item 
Sections~\ref{sec:mpit:cvar}, \ref{sec:mpit:pvar}, and~\ref{sec:mpit:cat}
on pages~\pageref{sec:mpit:cvar}, \pageref{sec:mpit:pvar}, and \pageref{sec:mpit:cat}.
\newline 
Clarified that \code{NULL} parameters can be provided
\mpifuncindex{MPI\_T\_CVAR\_GET\_INFO}%
\mpifuncindex{MPI\_T\_PVAR\_GET\_INFO}%
\mpifuncindex{MPI\_T\_CATEGORY\_GET\_INFO}%
in\\
\mpiskipfunc{MPI\_T\_\{CVAR$|$PVAR$|$CATEGORY\}\_GET\_INFO} routines. 

% 06.--- MPI-3.1 Ticket 377 and 400
\item 
Sections~\ref{sec:mpit:cvar}, \ref{sec:mpit:pvar}, \ref{sec:mpit:cat}, and~\ref{sec:mpit:retcodes}
on pages~\pageref{sec:mpit:cvar}, \pageref{sec:mpit:pvar}, \pageref{sec:mpit:cat}, and~\pageref{sec:mpit:retcodes}.
\newline 
New routines \mpifunc{MPI\_T\_CVAR\_GET\_INDEX}, \mpifunc{MPI\_T\_PVAR\_GET\_INDEX},
\mpifunc{MPI\_T\_CATEGORY\_GET\_INDEX}, were added to support retrieving
indices of variables and categories. 
The error codes \error{MPI\_T\_ERR\_INVALID} and \error{MPI\_T\_ERR\_INVALID\_NAME}
were added to indicate invalid uses of the interface.

\end{enumerate}


\section{Changes from Version 2.2 to Version 3.0}
\label{subsec:22to30}
 
% The following MPI-3.0 tickets are handled in a special way: 
%   #228 because MPIR is a stand-alone-document 
% 
% The following tickets do not have a change-log entry:
%   #158 #170 #172 #179 #181 #182 #185 #186 #187 #191 #201 #207 #215 #222
%   #229-A #251-(helper) #240-L #241-M #246-R #249-U #253-X
%   #254 #255 #259 #262 #263 #272 #278 #293 #298 
%   #308 #309 #312 #316 #317 #321 #329 #330 #331 #332 #333 #334 #335 #341
% 
% Additional correction in closed tickets:
%   #183 #197 #216 #260 #267 #268
%
% The following tickets did not pass MPI-2.2 final votes
%
% The following tickets are MPI-2.2-errata tickets
% (which should be mentioned below): 
%   #166 Item 1, #171 #192 #202 #340
% 
% Tickets that are expected to go to MPI-next (not yet moved):
%   #32 #47 #59 #165 #176 #195 #198 #310 #311 #319
%   #273 (under discussion)
%   #323 #326 #327 (FT)
%
% Tickets that are expected to be withdrawn (not yet withdrawn):
%   #102 #177 #178 #190 #315
%
% Tickets that need a change-log and do not yet have one and passed:
%   #38 #125 #126 #140 #162 #168, #171 2.2-ERRATA, 
%   #204 #219 #258 #265 #266 #270 #274 #280 #284 #286 #287 #299 #305 #322   
%
% Tickets that need a change-log and do not yet have one and passed in last meeting:
%   #166 Item 1 2.2-ERRATA, #192 2.2-ERRATA, #202 2.2-ERRATA, 
%   #256 #271 #281 #294 #300 #303 #313  #328
% 
% Tickets with change-log done: 
%   #109, #229ff (Fortran), #279 (obsolete due to #281)

\subsection{Fixes to Errata in Previous Versions of \texorpdfstring{\MPI/}{MPI}} 
\label{sec:changes-22-30-errata}
 
\begin{enumerate}

% 01.--- MPI-3.0 Ticket 343 = MPI-2.2 ERRATUM
\item 
Sections~\ref{sec:fortran-binding-issues} and~\ref{sec:c-binding-issues} 
on pages~\pageref{sec:fortran-binding-issues} and~\pageref{sec:c-binding-issues}, and
\MPIIIDOTII/ Section~2.6.2 on page~17, lines~41--42, 
Section~2.6.3 on page~18, lines~15--16, and 
Section~2.6.4 on page~18, lines~40--41.
\newline 
This is an \MPIII/ erratum: 
The scope for the reserved prefix \code{MPI\_} and 
the C++ namespace \code{MPI} is now any name as originally intended in \MPII/. 

% 02.--- MPI-3.0 Ticket 340 =  MPI-2.2 ERRATUM
\item 
Sections~\ref{subsec:pt2pt-messagedata},
\ref{coll-predefined-op},
\ref{subsec:ext32} Table~\ref{table:io:extsizes},
%-REMOVED-C++ \ref{sec:c++datatypes} Table \ref{tab:cpp-basic-datatypes},
and Annex~\ref{subsec:annexa-const} 
on pages~\pageref{subsec:pt2pt-messagedata},
\pageref{coll-predefined-op},
\pageref{table:io:extsizes}, 
%-REMOVED-C++ \pageref{tab:cpp-basic-datatypes},
and \pageref{subsec:annexa-const}, and
\MPIIIDOTII/ Sections 3.2.2, 5.9.2, 13.5.2 Table 13.2, 16.1.16 Table 16.1, and Annex A.1.1 on pages 27, 164, 433, 472 and 513
\newline 
This is an \MPIIIDOTII/ erratum: 
New named predefined datatypes \type{MPI\_CXX\_BOOL}, \type{MPI\_CXX\_FLOAT\_COMPLEX}, 
\type{MPI\_CXX\_DOUBLE\_COMPLEX}, and \type{MPI\_CXX\_LONG\_DOUBLE\_COMPLEX} were added
in C and Fortran corresponding to the C++ types 
\ctype{bool}, \ctype{std::complex<float>}, \ctype{std::complex<double>}, and \ctype{std::complex<long double>}. 
These datatypes also correspond to the deprecated C++ predefined datatypes 
\mpidtypeskip{MPI::BOOL}, \mpidtypeskip{MPI::COMPLEX}, \mpidtypeskip{MPI::DOUBLE\_COMPLEX}, and \mpidtypeskip{MPI::LONG\_DOUBLE\_COMPLEX},
which were removed in \MPIIIIDOTO/. 
The nonstandard C++ types \ctype{Complex<...>} were substituted by the standard types \ctype{std::complex<...>}.  

% 03.--- MPI-3.0 Ticket 340 = MPI-2.2 ERRATUM
\item 
Sections~\ref{coll-predefined-op} on pages~\pageref{coll-predefined-op} and
\MPIIIDOTII/ Section 5.9.2, page 165, line 47.
\newline 
This is an \MPIIIDOTII/ erratum: 
\type{MPI\_C\_COMPLEX} was added to the ``Complex'' reduction group. 

% 04.--- MPI-3.0 Ticket 192 = MPI-2.2 ERRATUM
\item 
Section~\ref{subsec:topol-inquiry} on page~\pageref{subsec:topol-inquiry}, and 
\MPIIIDOTII/, Section 7.5.5 on page 257, C++ interface on page 264, line 3. 
\newline 
This is an \MPIIIDOTII/ erratum: 
The argument \mpiarg{rank} was removed and \mpiarg{in/outdegree} are 
now defined as \code{int\& indegree} and \code{int\& outdegree}
in the C++ interface of 
\mpifunc{MPI\_DIST\_GRAPH\_NEIGHBORS\_COUNT}. 

% 05.--- MPI-3.0 Ticket 171 = MPI-2.2 ERRATUM
\item 
Section~\ref{subsec:ext32}, Table~\ref{table:io:extsizes} on page~\pageref{table:io:extsizes}, and
\MPIIIDOTII/, Section 13.5.3, Table 13.2 on page 433.
\newline 
This was an \MPIIIDOTII/ erratum: 
The \type{MPI\_C\_BOOL} \mpistring{external32} representation is corrected to a 1-byte size.  

% 06.--- MPI-3.0 Ticket 166 Item 1 = MPI-2.2 ERRATUM
\item 
{\raggedright
%-REMOVED-C++ Section~\ref{sec:c++datatypes} on page~\pageref{sec:c++datatypes}, and
%-REMOVED-C++ \newline 
\MPIIIDOTII/ Section 16.1.16 on page 471, line 45.
\newline 
This is an \MPIIIDOTII/ erratum: 
The constant \mpidtypeskip{MPI::\_LONG\_LONG} should be \mpidtypeskip{MPI::LONG\_LONG}.

}
% 07.--- MPI-3.0 Ticket 202 = MPI-2.2 ERRATUM
\item 
Annex~\ref{subsec:annexa-const} on page~\pageref{subsec:annexa-const}, Table ``Optional datatypes (Fortran),'' and
\newline 
\MPIIIDOTII/, Annex A.1.1, Table on page 517, lines 34, and 37--41. 
\newline 
This is an \MPIIIDOTII/ erratum: 
The C++ datatype handles \mpidtypeskip{MPI::INTEGER16}, \mpidtypeskip{MPI::REAL16}, 
\mpidtypeskip{MPI::F\_COMPLEX4}, \mpidtypeskip{MPI::F\_COMPLEX8}, \mpidtypeskip{MPI::F\_COMPLEX16}, \mpidtypeskip{MPI::F\_COMPLEX32}
were added to the table. 
 
\end{enumerate}

\subsection{Changes in \texorpdfstring{\MPIIIIDOTO/}{MPI-3.0}}
 
\begin{enumerate}

% 01.--- MPI-3.0 Ticket 281 
\item 
Section~\ref{sec:deprecated} on page~\pageref{sec:deprecated},
Section~\ref{sec:rm-cpp} on page~\pageref{sec:rm-cpp} and all other chapters. 
\newline 
The C++ bindings were removed from the standard.
See errata 
in Section~\ref{sec:changes-22-30-errata} on page~\pageref{sec:changes-22-30-errata}
for the latest changes to the \MPI/ C++ binding defined in \MPIIIDOTII/. 
\newline This change may affect backward compatibility.

% 02.--- MPI-3.0 Ticket 303 
\item
{\raggedright
Section~\ref{sec:deprecated} on page~\pageref{sec:deprecated},
Section~\ref{sec:deprecated:since20} on page~\pageref{sec:deprecated:since20} and
Section~\ref{sec:rm-mpii} on page~\pageref{sec:rm-mpii}.
\newline 
The deprecated functions 
\mpifunc{MPI\_TYPE\_HVECTOR},
\mpifunc{MPI\_TYPE\_HINDEXED},
\mpifunc{MPI\_TYPE\_STRUCT},
\mpifunc{MPI\_ADDRESS},
\mpifunc{MPI\_TYPE\_EXTENT},
\mpifunc{MPI\_TYPE\_LB},
\mpifunc{MPI\_TYPE\_UB},
\mpifunc{MPI\_ERRHANDLER\_CREATE} 
(and its callback function prototype \mpicallback{MPI\_Handler\_function}),
\mpifunc{MPI\_ERRHANDLER\_SET},
\mpifunc{MPI\_ERRHANDLER\_GET},
the deprecated special datatype handles
\mpidtype{MPI\_LB},
\mpidtype{MPI\_UB},
and the constants
\mpiconst{MPI\_COMBINER\_HINDEXED\_INTEGER},
\mpiconst{MPI\_COMBINER\_HVECTOR\_INTEGER},
\mpiconst{MPI\_COMBINER\_STRUCT\_INTEGER}
were removed from the standard. 
\newline This change may affect backward compatibility.

}
% 03.--- MPI-3.0 Ticket 140 
\item 
Section~\ref{terms:procedure-specification} on page~\pageref{terms:procedure-specification}. 
\newline 
Clarified parameter usage for \IN\ parameters. C bindings are now const-correct where
backward compatibility is preserved. 

% 04.--- MPI-3.0 Ticket 294 
\item 
Section~\ref{subsec:named-constants} on page~\pageref{subsec:named-constants} and
Section~\ref{subsec:topol-distgraph-constructor} on page~\pageref{subsec:topol-distgraph-constructor}. 
\newline 
The recommended C implementation value for \mpiconst{MPI\_UNWEIGHTED} 
changed from NULL to non-NULL. 
An additional weight array constant (\mpiconst{MPI\_WEIGHTS\_EMPTY}) was introduced.

% 05.--- MPI-3.0 Ticket 204 
\item 
Section~\ref{subsec:named-constants} on page~\pageref{subsec:named-constants} and 
Section~\ref{subsec:inquiry-version} on page~\pageref{subsec:inquiry-version}. 
\newline 
Added the new routine \mpifunc{MPI\_GET\_LIBRARY\_VERSION} 
to query library specific versions,
and the new constant \mpiconst{MPI\_MAX\_LIBRARY\_VERSION\_STRING}.

% 06.--- MPI-3.0 Ticket 265, 1st entry 
\item 
Sections~\ref{subsec:count}, \ref{subsec:pt2pt-messagedata},
\ref{table:pttopt:datatypes:c_f}, \ref{coll-predefined-op}, 
on pages~\pageref{subsec:count}, \pageref{subsec:pt2pt-messagedata},
\pageref{table:pttopt:datatypes:c_f}, \pageref{coll-predefined-op},
Sections~\ref{sec:pt2pt-datatype}, % MPI_TYPE_SIZE_X
\ref{subsec:pt2pt-extent}, % MPI_TYPE_GET_EXTENT_X
\ref{subsec:pt2pt-true-extent}, % MPI_TYPE_GET_TRUE_EXTENT_X
\ref{subsec:pt2pt-datatypeuse}, % MPI_GET_EXTENTS_X
\ref{func:mpi-status-set-elements-x} % MPI_STATUS_SET_ELEMENTS_X
on pages~\pageref{sec:pt2pt-datatype}, % MPI_TYPE_SIZE_X
\pageref{subsec:pt2pt-extent}, % MPI_TYPE_GET_EXTENT_X
\pageref{subsec:pt2pt-true-extent}, % MPI_TYPE_GET_TRUE_EXTENT_X
\pageref{subsec:pt2pt-datatypeuse}, % MPI_GET_EXTENTS_X
\pageref{func:mpi-status-set-elements-x}, % MPI_STATUS_SET_ELEMENTS_X
and
Annex~\ref{subsec:annexa-const} on page~\pageref{subsec:annexa-const}.
\newline 
New inquiry functions, \mpifunc{MPI\_TYPE\_SIZE\_X}, \mpifunc{MPI\_TYPE\_GET\_EXTENT\_X},
\mpifunc{MPI\_TYPE\_GET\_TRUE\_EXTENT\_X}, and \mpifunc{MPI\_GET\_ELEMENTS\_X}, return their results
as an \mpictype{MPI\_Count} value, which is a new type large enough to represent element
counts in memory, file views, etc. A new function, \mpifunc{MPI\_STATUS\_SET\_ELEMENTS\_X},
modifies the opaque part of an \mpictype{MPI\_Status} object so that a call to \mpifunc{MPI\_GET\_ELEMENTS\_X}
returns the provided \mpictype{MPI\_Count} value (in Fortran, \mpiftypekind{COUNT}).
The corresponding predefined datatype is \mpidtype{MPI\_COUNT}.  

% 07.--- MPI-3.0 Ticket 125+126
\item 
Chapter~\ref{chap:pt2pt} on page~\pageref{chap:pt2pt} through
Chapter~\ref{chap:binding-2} on page~\pageref{chap:binding-2}.
\newline 
In the C language bindings, the array-arguments' interfaces were modified to consistently 
use \code{{[}{]}} instead of \code{*}.

Exceptions are \mpifunc{MPI\_INIT}, which continues to use 
\code{char ***argv} 
(correct because of subtle rules regarding the use of 
the \code{\&} operator with \code{char *argv[]}), 
and \mpifunc{MPI\_INIT\_THREAD}, which is changed to be 
consistent with \mpifunc{MPI\_INIT}.
 
% 08.--- MPI-3.0 Ticket 265, 2nd entry 
\item 
Sections~\ref{subsec:pt2pt-status}, % MPI_GET_COUNT
\ref{subsec:pt2pt-addfunc}, % MPI_TYPE_SIZE
\ref{subsec:pt2pt-datatypeuse}, % MPI_GET_ELEMENTS
\ref{sec:pt2pt-packing} % MPI_PACK_SIZE
on pages~\pageref{subsec:pt2pt-status}, 
\pageref{subsec:pt2pt-addfunc},
\pageref{subsec:pt2pt-datatypeuse}, 
\pageref{sec:pt2pt-packing}.
\newline 
The functions \mpifunc{MPI\_GET\_COUNT} and \mpifunc{MPI\_GET\_ELEMENTS} were
defined to set the \mpiarg{count} argument to \mpiconst{MPI\_UNDEFINED} when that argument
would overflow. 
The functions \mpifunc{MPI\_PACK\_SIZE} and \mpifunc{MPI\_TYPE\_SIZE} were 
defined to set the \mpiarg{size} argument to \mpiconst{MPI\_UNDEFINED} when that argument 
would overflow. In all other \MPIIIDOTII/
routines, the type and semantics of the count arguments remain unchanged, i.e., \ctype{int}
or \ftype{INTEGER}. 
 
% 09.--- MPI-3.0 Ticket 229.2, 229.4 (and MPROBE ticket 38+274)
\item 
Section~\ref{sec:pt2pt-status-ignore} on page~\pageref{sec:pt2pt-status-ignore}, and
Section~\ref{sec:pt2pt-probe} on page~\pageref{sec:pt2pt-probe}.
\newline 
\mpiconst{MPI\_STATUS\_IGNORE} can also be used in 
\mpifunc{MPI\_IPROBE}, \mpifunc{MPI\_PROBE},
\mpifunc{MPI\_IMPROBE}, and \mpifunc{MPI\_MPROBE}.

% 10.--- MPI-3.0 Ticket 256+328 
\item 
Section~\ref{sec:pt2pt-probe} on page~\pageref{sec:pt2pt-probe} and  
Section~\ref{sec:pt2pt-nullproc} on page~\pageref{sec:pt2pt-nullproc}. 
\newline 
The use of \mpiconst{MPI\_PROC\_NULL} in probe operations was clarified. 
A special predefined message \mpiconst{MPI\_MESSAGE\_NO\_PROC} 
was defined for the use of matching probe 
(i.e., the new \mpifunc{MPI\_MPROBE} and \mpifunc{MPI\_IMPROBE}) 
with \mpiconst{MPI\_PROC\_NULL}. 
 
% 11.--- MPI-3.0 Tickets 38, 274
\item 
Sections~\ref{sec:matching-probe}, 
\ref{sec:matched-receive}, 
\ref{sec:misc-handleconvert}, 
\ref{subsec:annexa-const}
on pages~\pageref{sec:matching-probe},
\pageref{sec:matched-receive},
\pageref{sec:misc-handleconvert},
\pageref{subsec:annexa-const}. 
\newline 
Like \mpifunc{MPI\_PROBE} and \mpifunc{MPI\_IPROBE}, the new \mpifunc{MPI\_MPROBE} and \mpifunc{MPI\_IMPROBE} 
operations allow incoming messages to be queried without actually receiving them, except that
\mpifunc{MPI\_MPROBE} and \mpifunc{MPI\_IMPROBE} provide a mechanism to receive the specific message
with the new routines \mpifunc{MPI\_MRECV} and \mpifunc{MPI\_IMRECV} regardless of other intervening probe or receive operations. 
The opaque object \mpictype{MPI\_Message}, the null handle \mpiconst{MPI\_MESSAGE\_NULL}, and
the conversion functions  
\cfunc{MPI\_Message\_c2f}\index{MPI\_MESSAGE\_C2F} and
\cfunc{MPI\_Message\_f2c}\index{MPI\_MESSAGE\_F2C} 
were defined. 

% 12.--- MPI-3.0 Ticket 280 
\item 
Section~\ref{subsec:pt2pt-datatypeconst} on page~\pageref{subsec:pt2pt-datatypeconst} and
Section~\ref{subsec:typeaccess} on page~\pageref{subsec:typeaccess}. 
\newline 
The routine \mpifunc{MPI\_TYPE\_CREATE\_HINDEXED\_BLOCK} 
and constant \mpiconst{MPI\_COMBINER\_HINDEXED\_BLOCK} were added.

% 13.--- MPI-3.0 Ticket 109 
\item 
Chapter~\ref{chap:coll} on page~\pageref{chap:coll} and
Section~\ref{sec:nbcoll} on page~\pageref{sec:nbcoll}. 
\newline 
Added nonblocking interfaces to all collective operations.

% 14.--- MPI-3.0 Ticket 271 
\item 
Sections~\ref{subsec:context-intracomconst}, 
\ref{sec:comm-info}, % 6.4.4 Communicator Info, in context.tex
\ref{subsec:window-info}, % 11.2.3 Window Info, in one-sided-2.tex 
on pages~\pageref{subsec:context-intracomconst},
\pageref{sec:comm-info}, % 6.4.4 Communicator Info
\pageref{subsec:window-info}. % 11.2.3 Window Info
\newline 
The new routines \mpifunc{MPI\_COMM\_DUP\_WITH\_INFO}, 
\mpifunc{MPI\_COMM\_SET\_INFO}, \mpifunc{MPI\_COMM\_GET\_INFO}, 
\mpifunc{MPI\_WIN\_SET\_INFO}, and \mpifunc{MPI\_WIN\_GET\_INFO} were\flushline % added to correct overful box into margin
added.
The routine \mpifunc{MPI\_COMM\_DUP} must also duplicate 
info hints. 

% 15.--- MPI-3.0 Ticket 168 
\item 
Section~\ref{subsec:context-intracomconst} on page~\pageref{subsec:context-intracomconst}. 
\newline 
Added \mpifunc{MPI\_COMM\_IDUP}. 

% 16.--- MPI-3.0 Ticket 286 
\item 
Section~\ref{subsec:context-intracomconst} on page~\pageref{subsec:context-intracomconst}. 
\newline 
Added the new communicator construction routine \mpifunc{MPI\_COMM\_CREATE\_GROUP}, which is
invoked only by the processes in the group of the new communicator being constructed.

% 17.--- MPI-3.0 Ticket 287 
\item 
Section~\ref{subsec:context-intracomconst} on page~\pageref{subsec:context-intracomconst}. 
\newline 
Added the \mpifunc{MPI\_COMM\_SPLIT\_TYPE} routine
and the communicator split type constant \mpiconst{MPI\_COMM\_TYPE\_SHARED}.

% 18.--- MPI-3.0 Ticket 305 
\item 
Section~\ref{subsec:context-intercomm} on page~\pageref{subsec:context-intercomm}. 
\newline 
In \MPIIIDOTII/, communication involved in an  \mpifunc{MPI\_INTERCOMM\_CREATE} operation 
could interfere with point-to-point communication on the parent 
communicator with the same tag or \mpiconst{MPI\_ANY\_TAG}.  This interference has 
been removed in \MPIIIIDOTO/.

% 19.--- MPI-3.0 Ticket 219 
\item 
Section~\ref{sec:ei-naming} on page~\pageref{sec:ei-naming}. 
\newline 
Section 6.8 on page 238.
The constant \mpiconst{MPI\_MAX\_OBJECT\_NAME} also applies for type and window names. 

% 20.--- MPI-3.0 Ticket 162 
\item 
Section~\ref{subsec:topol-lowlevel} on page~\pageref{subsec:topol-lowlevel}. 
\newline 
\mpifunc{MPI\_CART\_MAP} can also be used for a zero-dimensional topologies. 

% 21.--- MPI-3.0 Ticket 258, 299 
\item 
Section~\ref{sec:sparsecoll} on page~\pageref{sec:sparsecoll} and
Section~\ref{sec:sparsecoll-nbc} on page~\pageref{sec:sparsecoll-nbc}.
\newline 
The following neighborhood collective communication routines 
were added to support sparse communication on virtual topology grids: 
\mpifunc{MPI\_NEIGHBOR\_ALLGATHER}, \mpifunc{MPI\_NEIGHBOR\_ALLGATHERV}, 
\mpifunc{MPI\_NEIGHBOR\_ALLTOALL}, \mpifunc{MPI\_NEIGHBOR\_ALLTOALLV},
\mpifunc{MPI\_NEIGHBOR\_ALLTOALLW} and the nonblocking variants  
\mpifunc{MPI\_INEIGHBOR\_ALLGATHER}, \mpifunc{MPI\_INEIGHBOR\_ALLGATHERV},
\mpifunc{MPI\_INEIGHBOR\_ALLTOALL}, \mpifunc{MPI\_INEIGHBOR\_ALLTOALLV}, and
\mpifunc{MPI\_INEIGHBOR\_ALLTOALLW}.
The displacement arguments in \mpifunc{MPI\_NEIGHBOR\_ALLTOALLW}
and \mpifunc{MPI\_INEIGHBOR\_ALLTOALLW} were defined as address size integers.
In \mpifunc{MPI\_DIST\_GRAPH\_NEIGHBORS}, an ordering rule was added
for communicators created with \mpifunc{MPI\_DIST\_GRAPH\_CREATE\_ADJACENT}.  

% 22.--- MPI-3.0 Ticket 313 
\item 
Section~\ref{sec:inquiry-startup} on page~\pageref{sec:inquiry-startup} and
Section~\ref{sec:ei-threads-initialization} on page~\pageref{sec:ei-threads-initialization}. 
\newline 
The use of \mpifunc{MPI\_INIT}, \mpifunc{MPI\_INIT\_THREAD} and \mpifunc{MPI\_FINALIZE}
was clarified. 
After \MPI/ is initialized, the application can access information about the execution
environment by querying the new predefined info object \mpiconst{MPI\_INFO\_ENV}.

% 23.--- MPI-3.0 Ticket 266, 1st part 
\item 
Section~\ref{sec:inquiry-startup} on page~\pageref{sec:inquiry-startup}. 
\newline 
Allow calls to \mpiskipfunc{MPI\_T} routines before \mpifunc{MPI\_INIT}
and after \mpifunc{MPI\_FINALIZE}. 

% 24.--- MPI-3.0 Tickets 270, 284, 300 
\item 
Chapter~\ref{chap:one-side-2} on page~\pageref{chap:one-side-2}.
\newline 
Substantial revision of the entire One-sided chapter, with new routines for window creation, 
additional synchronization methods in passive target communication, new one-sided communication routines, 
a new memory model, and other changes. 

% 25.--- MPI-3.0 Ticket 266, 2nd part 
\item 
Section~\ref{sec:mpit} on page~\pageref{sec:mpit}. 
\newline 
A new \MPI/ Tool Information Interface was added.
 
\noindent
The following changes are related to the Fortran language support.
 
% 26.--- MPI-3.0 Ticket 230-B
% --- MPI-3.0 Ticket 247-S
% --- MPI-3.0 Ticket 248-T
\item 
Section~\ref{terms:procedure-specification} on page~\pageref{terms:procedure-specification}, 
and Sections~\ref{f90:overview}, \ref{f90:mpif08}, \ref{sec:f90:requirements} 
on pages~\pageref{f90:overview}, \pageref{f90:mpif08}, and \pageref{sec:f90:requirements}.
\newline 
The new \code{mpi\_f08} Fortran module was introduced. 
 
% 27.--- MPI-3.0 Ticket 231-C
\item 
Section~\ref{terms:opaque-objects} on page~\pageref{terms:opaque-objects}, 
and Sections~\ref{f90:mpif08}, \ref{f90:extended}, \ref{sec:f90:requirements} 
on pages~\pageref{f90:mpif08}, \pageref{f90:extended}, and~\pageref{sec:f90:requirements}.
\newline 
Handles to opaque objects were defined as named types within the \code{mpi\_f08} Fortran module. 
The operators \code{.EQ.}, \code{.NE.}, \code{==}, and \code{/=} were overloaded to allow the comparison of these handles. 
The handle types and the overloaded operators are also available through the \code{mpi} Fortran module.
 
%-Rab %
%-Rab % Ticket 244-P was removed by decision at Chicago meeting, Jul 19, 2011
%-Rab % Question:    Overloading instead of IGNORE constants?  
%-Rab % Straw-votes: 1 x Yes, 11 x No, 8 x Abstain 
%-Rab % 
%-Rab % --- MPI-3.0 Ticket 244-P
%-Rab \item 
%-Rab {%
%-Rab Section~\ref{subsec:array-arguments} on page~\pageref{subsec:array-arguments},
%-Rab Section~\ref{sec:pt2pt-status-ignore} on page~\pageref{sec:pt2pt-status-ignore}, 
%-Rab Section~\ref{subsec:topol-distgraph-constructor} on page~\pageref{subsec:topol-distgraph-constructor},
%-Rab Section~\ref{subsec:topol-inquiry} on page~\pageref{subsec:topol-inquiry},
%-Rab Section~\ref{sec:processes} on page~\pageref{sec:processes},
%-Rab Section~\ref{subsec:multiple} on page~\pageref{subsec:multiple}, 
%-Rab Section~\ref{sec:ei-gr} on page~\pageref{sec:ei-gr}, and 
%-Rab Section~\ref{f90:mpif08} on page~\pageref{f90:mpif08}.
%-Rab \newline 
%-Rab With the \code{mpi\_f08} module, optional arguments through function overloading
%-Rab are used instead of 
%-Rab \mpiconst{MPI\_STATUS\_IGNORE}, \mpiconst{MPI\_STATUSES\_IGNORE},
%-Rab \mpiconst{MPI\_ERRCODES\_IGNORE},
%-Rab and \mpiconst{MPI\_UNWEIGHTED}.
%-Rab }% 
 
% 28.--- MPI-3.0 Ticket 234-F
% --- MPI-3.0 Ticket 234-G
% --- MPI-3.0 Ticket 236-H
\item 
Sections~\ref{subsec:named-constants}, \ref{sub:choice} 
on pages~\pageref{subsec:named-constants}, \pageref{sub:choice},  
Sections~\ref{f90:overview}, \ref{sec:misc-problems}, 
\ref{sec:f90-problems:strong-typing}, \ref{sec:misc-sequence}, 
\ref{sec:f90-problems:vector-subscripts}
on pages~\pageref{f90:overview}, \pageref{sec:misc-problems}, 
\pageref{sec:f90-problems:strong-typing}, \pageref{sec:misc-sequence}, 
\pageref{sec:f90-problems:vector-subscripts}, and
Sections~\ref{f90:mpif08}, \ref{f90:extended}, \ref{sec:f90:requirements}
on pages~\pageref{f90:mpif08}, \pageref{f90:extended}, \pageref{sec:f90:requirements}.
\newline 
Within the \code{mpi\_f08} Fortran module, 
choice buffers were defined as assumed-type and assumed-rank 
according to Fortran 2008 with TS 29113\mpitermindex{TS 29113}
\cite{fortran2008:ts29113-interop},
and the compile-time constant \mpiconst{MPI\_SUBARRAYS\_SUPPORTED} was set to \code{.TRUE.}.
With this, Fortran subscript triplets can be used in nonblocking \MPI/ operations;
vector subscripts are not supported in nonblocking operations. 
If the compiler does not support this Fortran
TS 29113\mpitermindex{TS 29113} feature,
the constant is set to \code{.FALSE.}.
 
% 29.--- MPI-3.0 Ticket 239-K
\item 
Section~\ref{sec:fortran-binding-issues} on page~\pageref{sec:fortran-binding-issues},
Section~\ref{f90:mpif08} on page~\pageref{f90:mpif08}, and
Section~\ref{sec:f90:requirements} on page~\pageref{sec:f90:requirements}.
\newline 
The \mpiarg{ierror} dummy arguments are \ftype{OPTIONAL} within the \code{mpi\_f08} Fortran module. 
 
% 30.--- MPI-3.0 Ticket 243-O
\item 
Section~\ref{subsec:pt2pt-status} on page~\pageref{subsec:pt2pt-status},
Sections~\ref{f90:mpif08}, \ref{f90:extended}, \ref{sec:f90:requirements}, 
on pages~\pageref{f90:mpif08}, \pageref{f90:extended}, \pageref{sec:f90:requirements}, 
and Section~\ref{sec:conversion:status} on page~\pageref{sec:conversion:status}. 
\newline 
Within the \code{mpi\_f08} Fortran module,
the status was defined as \mpiftype{TYPE(MPI\_Status)}\cdeclindex{MPI\_Status}. 
Additionally, within both the \code{mpi} and the \code{mpi\_f08} modules,
the constants \mpiconst{MPI\_STATUS\_SIZE}, 
\mpiconst{MPI\_SOURCE}, \mpiconst{MPI\_TAG}, \mpiconst{MPI\_ERROR},
and \mpiftype{TYPE(MPI\_Status)}\cdeclindex{MPI\_Status}
are defined.
New conversion routines were added: 
\mpifunc{MPI\_STATUS\_F2F08}, 
\mpifunc{MPI\_STATUS\_F082F}, 
\cfunc{MPI\_Status\_c2f08}\index{MPI\_STATUS\_C2F08}, and
\cfunc{MPI\_Status\_f082c}\index{MPI\_STATUS\_F082C}, 
In \code{mpi.h}, the new type \mpictype{MPI\_F08\_status}, 
and the external variables \mpiconst{MPI\_F08\_STATUS\_IGNORE} and \mpiconst{MPI\_F08\_STATUSES\_IGNORE} were added. 
 
% 31.--- MPI-3.0 Ticket 229.2 (and MPROBE ticket 38+274)
\item 
Section~\ref{sec:pt2pt-buffer} on page~\pageref{sec:pt2pt-buffer}.
\newline 
In Fortran with the \code{mpi} module or \code{mpif.h}, 
the type of the \mpiarg{buffer\_addr} argument of \mpifunc{MPI\_BUFFER\_DETACH}
is incorrectly defined and the argument is therefore unused.
 
% 32.--- MPI-3.0 Ticket 229.2 and 237-I
\item 
Section~\ref{sec:pt2pt-datatype} on page~\pageref{sec:pt2pt-datatype},
Section~\ref{subsec:pt2pt-markers} on page~\pageref{subsec:pt2pt-markers}, and
Section~\ref{sec:f90-problems:derived-types} on page~\pageref{sec:f90-problems:derived-types}.
\newline 
The Fortran alignments of basic datatypes within Fortran derived types
are implementation dependent; 
therefore it is recommended to use the \ftype{BIND(C)} attribute for derived types
in \MPI/ communication buffers. 
If an array of structures (in C/C++) or derived types (in Fortran) 
is to be used in \MPI/ communication buffers, 
it is recommended that the user creates a portable datatype handle and 
additionally applies \mpifunc{MPI\_TYPE\_CREATE\_RESIZED} to this datatype handle.

% 33.--- MPI-3.0 Ticket 252-W
% --- MPI-3.0 Ticket 251-V
\item 
Sections~%
\ref{sec:ei-dupdata},
\ref{subsec:coll-user-ops},
\ref{subsec:coll-process-local-reduction},
\ref{subsec:caching:datatypes},
\ref{sec:ei-naming},
\ref{subsec:inquiry-errhdlr-comm},
\ref{subsec:inquiry-errhdlr-windows},
\ref{subsec:inquiry-errhdlr-files},
\ref{sec:deprecated:since20},
\ref{f90:types}
on pages~%
\pageref{sec:ei-dupdata},
\pageref{subsec:coll-user-ops},
\pageref{subsec:coll-process-local-reduction},
\pageref{subsec:caching:datatypes},
\pageref{sec:ei-naming},
\pageref{subsec:inquiry-errhdlr-comm},
\pageref{subsec:inquiry-errhdlr-windows},
\pageref{subsec:inquiry-errhdlr-files},
\pageref{sec:deprecated:since20}, and
\pageref{f90:types}.
In some routines, the dummy argument names were changed because they
were identical to the Fortran keywords \ftype{TYPE} and \ftype{FUNCTION}.
The new dummy argument names must be used because the
\code{mpi} and \code{mpi\_f08} modules guarantee keyword-based actual argument lists. 
The argument name \mpiarg{type} was changed 
  % into \mpiarg{oldtype} 
  in \mpifunc{MPI\_TYPE\_DUP},
  % and into \mpiarg{datatype} in 
    the Fortran \ftype{USER\_FUNCTION} of \mpifunc{MPI\_OP\_CREATE},
    % and in 
    \mpifunc{MPI\_TYPE\_SET\_ATTR},
    \mpifunc{MPI\_TYPE\_GET\_ATTR},
    \mpifunc{MPI\_TYPE\_DELETE\_ATTR},
    \mpifunc{MPI\_TYPE\_SET\_NAME},
    \mpifunc{MPI\_TYPE\_GET\_NAME},
    \mpifunc{MPI\_TYPE\_MATCH\_SIZE},
  the callback prototype definition  
    \mpicallback{MPI\_Type\_delete\_attr\_function},
  and the predefined callback function
% RAB: The predefined callbacks are "constant" "function" pointers, therefore in both indexes
% WDG: Use mpiconstfunc for these entries; this ensures consistent indexing
    \mpiconstfunc{MPI\_TYPE\_NULL\_DELETE\_FN};
\mpiarg{function} was changed
  % into \mpiarg{user\_fn} 
    in \mpifunc{MPI\_OP\_CREATE},
  % into \mpiarg{comm\_errhandler\_fn} in
       \mpifunc{MPI\_COMM\_CREATE\_ERRHANDLER},
  % into \mpiarg{win\_errhandler\_fn} in
       \mpifunc{MPI\_WIN\_CREATE\_ERRHANDLER},
  % into \mpiarg{file\_errhandler\_fn} in
       \mpifunc{MPI\_FILE\_CREATE\_ERRHANDLER}, and
  % into \mpiarg{handler\_fn} in
       \mpifunc{MPI\_ERRHANDLER\_CREATE}.
For consistency reasons, \mpiarg{INOUBUF} was changed to \mpiarg{INOUTBUF}
  in \mpifunc{MPI\_REDUCE\_LOCAL},
and \mpiarg{intracomm} to \mpiarg{newintracomm} 
  in \mpifunc{MPI\_INTERCOMM\_MERGE}.

% 34.--- MPI-3.0 Ticket 322 
\item 
Section~\ref{subsec:context-cachecomm} on page~\pageref{subsec:context-cachecomm}. 
\newline 
It was clarified that in Fortran, the flag values 
returned by a \mpiarg{comm\_copy\_attr\_fn} callback,
% RAB: The predefined callbacks are "constant" "function" pointers, therefore in both indexes
including \mpiconstfunc{MPI\_COMM\_NULL\_COPY\_FN}
and \mpiconstfunc{MPI\_COMM\_DUP\_FN}, 
are \code{.FALSE.} and \code{.TRUE.}; see \mpifunc{MPI\_COMM\_CREATE\_KEYVAL}.

% 35.--- MPI-3.0 Ticket 245-Q
\item 
Section~\ref{sec:misc-memalloc} on page~\pageref{sec:misc-memalloc}.
\newline 
With the \code{mpi} and \code{mpi\_f08} Fortran modules, 
\mpifunc{MPI\_ALLOC\_MEM} now also supports \ftype{TYPE(C\_PTR)} C-pointers 
instead of only returning an address-sized integer that may be usable together 
with a nonstandard Cray-pointer.
 
% 36.--- MPI-3.0 Ticket 237-I
\item 
Section~\ref{sec:f90-problems:derived-types} on page~\pageref{sec:f90-problems:derived-types}, and
Section~\ref{sec:f90:requirements} on page~\pageref{sec:f90:requirements}.
\newline 
Fortran \ftype{SEQUENCE} and \ftype{BIND(C)} derived application types
can now be used as buffers in \MPI/ operations.
 
% 37.--- MPI-3.0 Ticket 238-J
\item 
Section~\ref{sec:misc-register} on page~\pageref{sec:misc-register} to 
Section~\ref{sec:f90-problems:perm-data-movements} on page~\pageref{sec:f90-problems:perm-data-movements}, 
Section~\ref{sec:f90:requirements} on page~\pageref{sec:f90:requirements}, and
Section~\ref{f90:syncreg} on page~\pageref{f90:syncreg}.
\newline 
The sections about Fortran optimization problems and their solutions were
partially rewritten and new methods are added, e.g., the use of the \ftype{ASYNCHRONOUS} attribute.
The constant \mpiconst{MPI\_ASYNC\_PROTECTS\_NONBLOCKING} tells whether
the semantics of the \ftype{ASYNCHRONOUS} attribute is extended to protect 
nonblocking operations.
The Fortran routine \mpifunc{MPI\_F\_SYNC\_REG} is added. 
\MPIIIIDOTO/ compliance for an \MPI/ library together with a Fortran compiler is
defined in Section~\ref{sec:f90:requirements}. 
 
% 38.--- MPI-3.0 Ticket 242-N
\item 
Section~\ref{f90:mpif08} on page~\pageref{f90:mpif08}.
\newline 
Within the \code{mpi\_f08} Fortran module,
dummy arguments are now declared with \ftype{INTENT=IN},
\ftype{OUT}, or \ftype{INOUT} as defined in the \code{mpi\_f08} interfaces.
 
% 39.--- MPI-3.0 Ticket 232-D
\item 
Section~\ref{f90:extended} on page~\pageref{f90:extended}, and
Section~\ref{sec:f90:requirements} on page~\pageref{sec:f90:requirements}.
\newline 
The existing \code{mpi} Fortran module must implement compile-time argument checking. 
 
% 40.--- MPI-3.0 Ticket 233-E
\item 
Section~\ref{f90:basic} on page~\pageref{f90:basic}.
\newline 
The use of the \code{mpif.h} Fortran include file is now strongly discouraged. 
 
% 41.--- MPI-3.0 Ticket 230-B
\item 
Section~\ref{subsec:annexa-const},
Table \textbf{Predefined functions} on page~\pageref{table:predefined-functions},
Section~\ref{subsec:annexa-prototypes} on page~\pageref{subsec:annexa-prototypes},
and Section~\ref{app:subsec:f08:context} on page~\pageref{app:subsec:f08:context}.
\newline 
Within the new \code{mpi\_f08} module, all callback prototype definitions are now
defined with explicit interfaces \ftype{PROCEDURE(MPI\_...)} that have the \ftype{BIND(C)} 
attribute; user-written callbacks must be modified if the \code{mpi\_f08} module is used.
 
% 42.--- MPI-3.0 Ticket 250-V
\item 
Section~\ref{subsec:annexa-prototypes} on page~\pageref{subsec:annexa-prototypes}.
\newline 
In some routines, the Fortran callback prototype names were changed 
from \mpiskipfunc{$\ldots$\_FN} to \mpiskipfunc{$\ldots$\_FUNCTION} to be 
consistent with the other language bindings.
 
\end{enumerate}


% --- MPI-2.2 Ticket 100
\section{Changes from Version 2.1 to Version 2.2}

% The following MPI-2.2 tickets are handled in a special way: 
%   #049 #061 #099 #100
%
% The following tickets do not have a change-log entry:
%   #001 #008 #009 #010 #011 #013 #016 #017 #021 #030 #037 #040 #041 #042 #044 #051
%   #053 #057 #060 #067 #070 #074 #077 #087 #089 #090 #091 #092 #093 #097
%   #103 #104 #105 #107 #113 #115 #116 #118 #120 #121 #122 #123 #124 #128 #132 #135
%   #136 #141 #142 #146 #148 #149
%
% The following tickets did not pass MPI-2.2 final votes 
%   With changelogs: #59 #80 #127  (removed with "%---" )

\begin{enumerate}

% --- MPI-2.2 Ticket 65
\item 
Section~\ref{subsec:named-constants} on page~\pageref{subsec:named-constants}. 
\newline 
It is now guaranteed that predefined named constant handles (as other constants) can be used in initialization expressions or assignments, i.e., also before the call to \mpifunc{MPI\_INIT}.
 
% --- MPI-2.2 Ticket 150
\item
Section~\ref{subsec:lang} on page~\pageref{subsec:lang}, % Sect. 2.6 p.15
%-REMOVED-WITH-MPI-3.0: Section~\ref{terms-cpp} on page~\pageref{terms-cpp}, and % Sect. 2.6.4 p.18
%-REMOVED-WITH-MPI-3.0:  Section~\ref{sec:cpp} on page~\pageref{sec:cpp}. % Sect. 16.1 p.449 newlabel
and Section~\ref{sec:rm-cpp} on page~\pageref{sec:rm-cpp}. % Sect. 16.1 p.449 newlabel
\newline
The C++ language bindings have been deprecated and may be removed in a future version of the \MPI/ specification.

% --- MPI-2.2 Ticket 63
\item
Section~\ref{subsec:pt2pt-messagedata} on page~\pageref{subsec:pt2pt-messagedata}. % Sect. 3.2.2 p.27
\newline
\mpidtype{MPI\_CHAR} for printable characters is now defined for C type char (instead of signed char). This change should not have any impact on applications nor on \MPI/ libraries (except some comment lines), because printable characters could and can be stored in any of the C types char, signed char, and unsigned char, and \mpidtype{MPI\_CHAR} is not allowed for predefined reduction operations.

% --- MPI-2.2 Ticket 18
\item
\typeindex{MPI\_INT8\_T}\typeindex{MPI\_INT16\_T}\typeindex{MPI\_INT32\_T}\typeindex{MPI\_INT64\_T}%
\typeindex{MPI\_UINT8\_T}\typeindex{MPI\_UINT16\_T}\typeindex{MPI\_UINT32\_T}\typeindex{MPI\_UINT64\_T}%
Section~\ref{subsec:pt2pt-messagedata} on page~\pageref{subsec:pt2pt-messagedata}. % Sect. 3.2.2 p.27
\newline
\constskip{MPI\_(U)INT\{8,16,32,64\}\_T}, 
\mpidtype{MPI\_AINT}, \mpidtype{MPI\_OFFSET}, \mpidtype{MPI\_C\_BOOL}, \mpidtype{MPI\_C\_COMPLEX}, \mpidtype{MPI\_C\_FLOAT\_COMPLEX}, \mpidtype{MPI\_C\_DOUBLE\_COMPLEX}, and \mpidtype{MPI\_C\_LONG\_DOUBLE\_COMPLEX} are now valid predefined \MPI/ datatypes. 

% --- MPI-2.2 Ticket 45+98
\item
Section~\ref{sec:pt2pt-modes} on page~\pageref{sec:pt2pt-modes}, % Sect. 3.4 p.38
Section~\ref{subsec:pt2pt-commstart} on page~\pageref{subsec:pt2pt-commstart}, % Sect. 3.7.2 p.49
Section~\ref{sec:pt2pt-persistent} on page~\pageref{sec:pt2pt-persistent}, and % Sect. 3.9 p.68
Section~\ref{sec:coll-intro} on page~\pageref{sec:coll-intro}. % Sect. 5.1 p.129 newlabel
\newline
The read access restriction on the send buffer for blocking, non blocking and collective API has been lifted. It is permitted to access for read the send buffer while the operation is in progress. 


% --- MPI-2.2 Ticket 50
\item
Section~\ref{sec:pt2pt-nonblock} on page~\pageref{sec:pt2pt-nonblock}. % Sect. 3.7 p.47
\newline
The Advice to users for IBSEND and IRSEND was slightly changed.

% --- MPI-2.2 Ticket 143
\item
Section~\ref{subsec:pt2pt-commend} on page~\pageref{subsec:pt2pt-commend}. % Sect. 3.7.3 p.52
\newline
The advice to free an active request was removed in the Advice to users for \mpifunc{MPI\_REQUEST\_FREE}.

% --- MPI-2.2 Ticket 137
\item
Section~\ref{subsec:pt2pt-teststatus} on page~\pageref{subsec:pt2pt-teststatus}. % Sect. 3.7.6 p.63 newlabel
\newline
\mpifunc{MPI\_REQUEST\_GET\_STATUS} changed to permit inactive or null requests as input.

% --- MPI-2.2 Ticket 31
\item
Section~\ref{sec:coll-alltoall} on page~\pageref{sec:coll-alltoall}. % Sect. 5.8 p.155
\newline
``In place'' option is added to \mpifunc{MPI\_ALLTOALL}, \mpifunc{MPI\_ALLTOALLV}, and \mpifunc{MPI\_ALLTOALLW} for intra-communicators.

% --- MPI-2.2 Ticket 64
\item
Section~\ref{coll-predefined-op} on page~\pageref{coll-predefined-op}. % Sect. 5.9.2 p.161
\newline
Predefined parameterized datatypes (e.g., returned by\\\mpifunc{MPI\_TYPE\_CREATE\_F90\_REAL}) and optional named predefined datatypes (e.g. \mpidtype{MPI\_REAL8}) have been added to the list of valid datatypes in reduction operations.

% --- MPI-2.2 Ticket 18
\item
\typeindex{MPI\_INT8\_T}\typeindex{MPI\_INT16\_T}\typeindex{MPI\_INT32\_T}\typeindex{MPI\_INT64\_T}%
\typeindex{MPI\_UINT8\_T}\typeindex{MPI\_UINT16\_T}\typeindex{MPI\_UINT32\_T}\typeindex{MPI\_UINT64\_T}%
Section~\ref{coll-predefined-op} on page~\pageref{coll-predefined-op}. % Sect. 5.9.2 p.161
\newline
\constskip{MPI\_(U)INT\{8,16,32,64\}\_T} are all considered C integer types for the purposes of the predefined reduction operators. \mpidtype{MPI\_AINT} and \mpidtype{MPI\_OFFSET} are considered Fortran integer types. \mpidtype{MPI\_C\_BOOL} is considered a Logical type. \mpidtype{MPI\_C\_COMPLEX}, \mpidtype{MPI\_C\_FLOAT\_COMPLEX}, \mpidtype{MPI\_C\_DOUBLE\_COMPLEX}, and \mpidtype{MPI\_C\_LONG\_DOUBLE\_COMPLEX} are considered Complex types.

% --- MPI-2.2 Ticket 24
\item
Section~\ref{subsec:coll-process-local-reduction} on page~\pageref{subsec:coll-process-local-reduction}. % Sect. 5.9.7 p.173 NEWsection
\newline
The local routines \mpifunc{MPI\_REDUCE\_LOCAL} and \mpifunc{MPI\_OP\_COMMUTATIVE} have been added. 

% --- MPI-2.2 Ticket 27
\item
Section~\ref{subsec:coll-reduce-scatter-block} on page~\pageref{subsec:coll-reduce-scatter-block}. % Sect. 5.10.1 p.174 NEWsection
\newline
The collective function \mpifunc{MPI\_REDUCE\_SCATTER\_BLOCK} is added to the \MPI/ standard.

% --- MPI-2.2 Ticket 94
\item
Section~\ref{coll-exscan} on page~\pageref{coll-exscan}. % Sect. 5.11.2 p.175 newlabel
\newline
Added in place argument to \mpifunc{MPI\_EXSCAN}.

% --- MPI-2.2 Ticket 19
\item
Section~\ref{subsec:context-intracomconst} on page~\pageref{subsec:context-intracomconst}, and % Sect. 6.4.2 p.194
Section~\ref{sec:context-intercom} on page~\pageref{sec:context-intercom}. % Sect. 6.6 p.209 newlabel
\newline
Implementations that did not implement \mpifunc{MPI\_COMM\_CREATE} on inter\hyphen/communicators will need to add that functionality. As the standard described the behavior of this operation on inter-communicators, it is believed that most implementations already provide this functionality. Note also that the C++ binding for both \mpifunc{MPI\_COMM\_CREATE} and \mpifunc{MPI\_COMM\_SPLIT} explicitly allow Intercomms.

% --- MPI-2.2 Ticket 66
\item
Section~\ref{subsec:context-intracomconst} on page~\pageref{subsec:context-intracomconst}. % Sect. 6.4.2 p.194
\newline
\mpifunc{MPI\_COMM\_CREATE} is extended to allow several disjoint subgroups as input if comm is an intra-communicator. If comm is an inter-communicator it was clarified that all processes in the same local group of comm must specify the same value for group.

% --- MPI-2.2 Ticket 33
\item
%-REMOVED-WITH-MPI-3.0: \cdeclindex{MPI::Distgraphcomm}%
Section~\ref{subsec:topol-distgraph-constructor} on page~\pageref{subsec:topol-distgraph-constructor}. % Sect. 7.5.3a p.247 NEWsection
\newline
New functions for a scalable distributed graph topology interface has been added. In this section, the functions \mpifunc{MPI\_DIST\_GRAPH\_CREATE\_ADJACENT} and \mpifunc{MPI\_DIST\_GRAPH\_CREATE}, the constants \mpiconst{MPI\_UNWEIGHTED}, and the derived C++ class Distgraphcomm were added.

% --- MPI-2.2 Ticket 33
\item
{\raggedright
Section~\ref{subsec:topol-inquiry} on page~\pageref{subsec:topol-inquiry}. % Sect. 7.5.4 p.248
\newline
For the scalable distributed graph topology interface, the 
functions \mpifunc{MPI\_DIST\_GRAPH\_NEIGHBORS\_COUNT} 
and \mpifunc{MPI\_DIST\_GRAPH\_NEIGHBORS} 
and the constant \mpiconst{MPI\_DIST\_GRAPH} were added.

}
% --- MPI-2.2 Ticket 3
\item
Section~\ref{subsec:topol-inquiry} on page~\pageref{subsec:topol-inquiry}. % Sect. 7.5.4 p.248
\newline
Remove ambiguity regarding duplicated neighbors with \mpifunc{MPI\_GRAPH\_NEIGHBORS} and \mpifunc{MPI\_GRAPH\_NEIGHBORS\_COUNT}. 

% --- MPI-2.2 Ticket 101
\item
Section~\ref{subsec:inquiry-version} on page~\pageref{subsec:inquiry-version}. % Sect. 8.1.1 p.259
\newline
The subversion number changed from 1 to 2.

\mpicallbackindex{MPI\_Comm\_errhandler\_fn}\mpicallbackindex{MPI\_Comm\_errhandler\_function}% 
\mpicallbackindex{MPI\_File\_errhandler\_fn}\mpicallbackindex{MPI\_File\_errhandler\_function}% 
\mpicallbackindex{MPI\_Win\_errhandler\_fn}\mpicallbackindex{MPI\_Win\_errhandler\_function}% 
% --- MPI-2.2 Ticket 7
\item
%corrected to
Section~\ref{sec:errorhandler} on page~\pageref{sec:errorhandler}, % Sect. 8.3 p.264
Section~\ref{sec:deprecated:since22} on page~\pageref{sec:deprecated:since22}, and % Sect. 15.2 p.448
Annex~\ref{subsec:annexa-prototypes} on page~\pageref{subsec:annexa-prototypes}. % Sect. A.1.3 p.500 newlabel
\newline
Changed function pointer typedef names \constskip{MPI\_\{Comm,File,Win\}\_errhandler\_fn} to \constskip{MPI\_\{Comm,File,Win\}\_errhandler\_function}. Deprecated old ``\_fn'' names. 


% --- MPI-2.2 Ticket 71
\item
Section~\ref{subsec:inquiry-startup-userfunc} on page~\pageref{subsec:inquiry-startup-userfunc}. % Sect. 8.7.1 p.283 newlabel
\newline
Attribute deletion callbacks on \mpiconst{MPI\_COMM\_SELF} are now called in LIFO order. Implementors must now also register all implementation-internal attribute deletion callbacks on \mpiconst{MPI\_COMM\_SELF} before returning from \mpifunc{MPI\_INIT}/\mpifunc{MPI\_INIT\_THREAD}.

% --- MPI-2.2 Ticket 43
\item
Section~\ref{sec:1sided-accumulate} on page~\pageref{sec:1sided-accumulate}. % Sect. 11.3.4 p.331
\newline
The restriction added in \MPI/ 2.1 that the operation \mpiconst{MPI\_REPLACE} in \mpifunc{MPI\_ACCUMULATE} can be used only with predefined datatypes has been removed. \mpiconst{MPI\_REPLACE} can now be used even with derived datatypes, as it was in \MPI/ 2.0. Also, a clarification has been made that \mpiconst{MPI\_REPLACE} can be used only in \mpifunc{MPI\_ACCUMULATE}, not in collective operations that do reductions, such as \mpifunc{MPI\_REDUCE} and others.

% --- MPI-2.2 Ticket 6
\item
\mpifuncindex{MPI\_GREQUEST\_START}% 
Section~\ref{sec:ei-gr} on page~\pageref{sec:ei-gr}. % Sect. 12.2 p.357
\newline
Add ``\code{*}'' to the \mpiarg{query\_fn}, \mpiarg{free\_fn}, and \mpiarg{cancel\_fn} arguments to the C++ binding for \mpiskipfunc{MPI::Grequest::Start()} for consistency with the rest of \MPI/ functions that take function pointer arguments. 

% --- MPI-2.2 Ticket 18
\item
{\raggedright
\typeindex{MPI\_INT8\_T}\typeindex{MPI\_INT16\_T}\typeindex{MPI\_INT32\_T}\typeindex{MPI\_INT64\_T}%
\typeindex{MPI\_UINT8\_T}\typeindex{MPI\_UINT16\_T}\typeindex{MPI\_UINT32\_T}\typeindex{MPI\_UINT64\_T}%
Section~\ref{subsec:ext32} on page~\pageref{subsec:ext32}, and  % Sect. 13.5.2 p.414
Table~\ref{table:io:extsizes} on page~\pageref{table:io:extsizes}. % Tab. 13.2 p.416
\newline
\constskip{MPI\_(U)INT\{8,16,32,64\}\_T}, \mpidtype{MPI\_AINT}, \mpidtype{MPI\_OFFSET}, \mpidtype{MPI\_C\_COMPLEX}, \mpidtype{MPI\_C\_FLOAT\_COMPLEX}, \mpidtype{MPI\_C\_DOUBLE\_COMPLEX},\\\mpidtype{MPI\_C\_LONG\_DOUBLE\_COMPLEX}, and \mpidtype{MPI\_C\_BOOL} are added as predefined datatypes in the
\mpistring{external32} representation.

}
% --- MPI-2.2 Ticket 55
\item
Section~\ref{sec:misc-attr} on page~\pageref{sec:misc-attr}. % Sect. 16.3.7 p.486
\newline
% {\bf !!!TODO!!! See Ticket -- proposed text:} 
The description was modified that it only describes how an \MPI/ implementation behaves, but not how \MPI/ stores attributes internally. The erroneous \MPIIIDOTI/ Example 16.17 was replaced with three new examples \ref{exa:misc-attr:c}, \ref{exa:misc-attr:f-old}, and \ref{exa:misc-attr:f-new} on pages~\pageref{exa:misc-attr:c:pagebegin}--\pageref{exa:misc-attr:f-new:pagebegin} explicitly detailing cross-language attribute behavior. Implementations that matched the behavior of the old example will need to be updated.  

% --- MPI-2.2 Ticket 4
\item
Annex~\ref{subsec:annexa-const} on page~\pageref{subsec:annexa-const}. % Sect. A.1.1 p.491
\newline
Removed type \constskip{MPI::Fint} (compare \mpictype{MPI\_Fint} in 
Section~\ref{subsec:annexa-type} on page~\pageref{subsec:annexa-type}). % Sect. A.1.2 p. 499 newlabel

% --- MPI-2.2 Ticket 18
\item
\typeindex{MPI\_INT8\_T}\typeindex{MPI\_INT16\_T}\typeindex{MPI\_INT32\_T}\typeindex{MPI\_INT64\_T}%
\typeindex{MPI\_UINT8\_T}\typeindex{MPI\_UINT16\_T}\typeindex{MPI\_UINT32\_T}\typeindex{MPI\_UINT64\_T}%
% Annex~\ref{sec:lang:const} on page~\pageref{sec:lang:const}, % Sect. A.1 p.491 
% It is better to use the following reference: 
Annex~\ref{subsec:annexa-const} on page~\pageref{subsec:annexa-const}. % Sect. A.1.1 p.491
Table \textbf{Named Predefined Datatypes}. % on p.494
\newline
Added \constskip{MPI\_(U)INT\{8,16,32,64\}\_T}, \mpidtype{MPI\_AINT}, \mpidtype{MPI\_OFFSET}, \mpidtype{MPI\_C\_BOOL}, \mpidtype{MPI\_C\_FLOAT\_COMPLEX}, \mpidtype{MPI\_C\_COMPLEX}, \mpidtype{MPI\_C\_DOUBLE\_COMPLEX}, and \mpidtype{MPI\_C\_LONG\_DOUBLE\_COMPLEX} are added as predefined datatypes.

  
\end{enumerate}

\section{Changes from Version 2.0 to Version 2.1} 
\label{sec:changes:last}
 

\begin{enumerate}

\item
Section~\ref{subsec:pt2pt-messagedata} on page~\pageref{subsec:pt2pt-messagedata},
%-REMOVED-WITH-MPI-3.0:                Section~\ref{sec:c++datatypes} on page~\pageref{sec:c++datatypes},
                                       and Annex~\ref{sec:lang:const} on page~\pageref{sec:lang:const}. 
% --- MPI-2.1, Ballot 1, Item 16 ---
\newline
In addition, the \mpidtype{MPI\_LONG\_LONG} should be added as an optional
type; it is a synonym for \mpidtypeskip{MPI\_LONG\_LONG\_INT}.

\item
Section~\ref{subsec:pt2pt-messagedata} on page~\pageref{subsec:pt2pt-messagedata},
%-REMOVED-WITH-MPI-3.0:                Section~\ref{sec:c++datatypes} on page~\pageref{sec:c++datatypes},
                                       and Annex~\ref{sec:lang:const} on page~\pageref{sec:lang:const}. 
% --- MPI-2.1, Review Item 15.h' ---
\newline
\mpidtype{MPI\_LONG\_LONG\_INT}, \mpidtype{MPI\_LONG\_LONG} (as synonym),\\\mpidtype{MPI\_UNSIGNED\_LONG\_LONG},   
\mpidtype{MPI\_SIGNED\_CHAR},
and \mpidtype{MPI\_WCHAR} are moved from optional to official and they are therefore
defined for all three language bindings.

\item
Section~\ref{subsec:pt2pt-status} on page~\pageref{subsec:pt2pt-status}.
% --- MPI-2.1, Ballot 2, Item 1 ---
\newline
\mpifunc{MPI\_GET\_COUNT} with zero-length datatypes:
The value returned as the \mpiarg{count} argument of
\mpifunc{MPI\_GET\_COUNT} for a datatype of length zero where zero bytes
have been transferred is zero.  If the number of bytes transferred is
greater than zero, \mpiconst{MPI\_UNDEFINED} is returned. 

\item
Section~\ref{sec:pt2pt-datatype} on page~\pageref{sec:pt2pt-datatype}.
% --- MPI-2.1, Ballot 4, Item 20 ---
\newline
General rule about derived datatypes:
Most datatype constructors have replication count or block length arguments.
Allowed values are 
{nonnegative}
integers. If the value is zero, no elements are
generated in the type map and there is no effect on datatype bounds or
extent. 

\item
Section~\ref{canonical_pack} on page~\pageref{canonical_pack}.
% --- MPI-2.1, Ballot 2, Item 4 ---
\newline
\mpidtype{MPI\_BYTE} should be used to send and receive data that is packed
using \mpifunc{MPI\_PACK\_EXTERNAL}.

\item
Section~\ref{subsec:coll-all-reduce} on page~\pageref{subsec:coll-all-reduce}.
% --- MPI-2.1, Ballot 2, Item 8 ---
\newline
If \mpiarg{comm} is an inter-communicator in \mpifunc{MPI\_ALLREDUCE}, then
both groups should
provide \mpiarg{count} and \mpiarg{datatype} arguments that specify the same type
signature
(i.e., it is not necessary that both groups provide the same \mpiarg{count} value). 

\item
Section~\ref{subsec:context-grpacc} on page~\pageref{subsec:context-grpacc}.
% --- MPI-2.1, Ballot 2, Item 2 ---
\newline
\mpifunc{MPI\_GROUP\_TRANSLATE\_RANKS} and \mpiconst{MPI\_PROC\_NULL}:
\mpiconst{MPI\_PROC\_NULL} is a valid rank for input to
\mpifunc{MPI\_GROUP\_TRANSLATE\_RANKS}, which returns
\mpiconst{MPI\_PROC\_NULL} as the translated rank.

\item
Section~\ref{sec:caching} on page~\pageref{sec:caching}.
% --- MPI-2.1, Ballot 2, Item 9 ---
\newline
About the attribute caching functions: 
\begin{implementors}
High-quality implementations should raise an error when a keyval
\mpifuncindex{MPI\_TYPE\_CREATE\_KEYVAL}%
\mpifuncindex{MPI\_COMM\_CREATE\_KEYVAL}%
\mpifuncindex{MPI\_WIN\_CREATE\_KEYVAL}%
that was created by a call to \mpiskipfunc{MPI\_\XXX/\_CREATE\_KEYVAL} is
used with an object of the wrong type with a call to
\mpiskipfunc{MPI\_YYY\_GET\_ATTR}, \mpiskipfunc{MPI\_YYY\_SET\_ATTR}, \mpiskipfunc{MPI\_YYY\_DELETE\_ATTR}, or
\mpiskipfunc{MPI\_YYY\_FREE\_KEYVAL}. To do so, it is necessary to maintain, with
each keyval, information on the type of the associated user
function.
\end{implementors}

\item
Section~\ref{sec:ei-naming} on page~\pageref{sec:ei-naming}.
% --- MPI-2.1, Ballot 4, Item 13.iii ---
\newline
In \mpifunc{MPI\_COMM\_GET\_NAME}: 
In C, a null character is additionally stored at \mpiarg{name{[}resultlen{]}}. 
\mpiarg{resultlen} cannot be larger then 
{\mpiconst{MPI\_MAX\_OBJECT\_NAME}}-1. 
In Fortran, name is padded on the right with blank characters. 
\mpiarg{resultlen} cannot be larger then 
{\mpiconst{MPI\_MAX\_OBJECT\_NAME}}. 

\item
Section~\ref{subsec:topol-overview} on page~\pageref{subsec:topol-overview}.
% --- MPI-2.1, Ballot 4, Item 11.c ---
\newline
About \mpifunc{MPI\_GRAPH\_CREATE} and \mpifunc{MPI\_CART\_CREATE}:
All input arguments must have identical values on all processes
of the group of \mpiarg{comm\_old}. 

\item
Section~\ref{subsec:topol-cartesian-constructor} on page~\pageref{subsec:topol-cartesian-constructor}.
% --- MPI-2.1, Ballot 4, Item 10.d ---
\newline
In \mpifunc{MPI\_CART\_CREATE}:
If \mpiarg{ndims} is zero then a zero-dimensional Cartesian topology is created. The
call is erroneous if it specifies a grid that is larger than the group size or if
\mpiarg{ndims} is negative.

\item
Section~\ref{subsec:topol-graph-constructor} on page~\pageref{subsec:topol-graph-constructor}.
% --- MPI-2.1, Ballot 4, Item 11.a ---
\newline
In \mpifunc{MPI\_GRAPH\_CREATE}:
If the graph is empty, i.e., \mpiarg{nnodes}\mpicode{ = 0}, 
then \mpiconst{MPI\_COMM\_NULL} is returned in all processes. 

\item
Section~\ref{subsec:topol-graph-constructor} on page~\pageref{subsec:topol-graph-constructor}.
% --- MPI-2.1, Ballot 4, Item 11.b ---
\newline
In \mpifunc{MPI\_GRAPH\_CREATE}:
A single process is allowed to be defined multiple times in the list of
neighbors of a process (i.e., there may be multiple edges between two
processes). A process is also allowed to be a neighbor to itself (i.e., a self
loop in the graph). The adjacency matrix is allowed to be nonsymmetric.
\begin{users} 
Performance implications of using multiple edges or a nonsymmetric
adjacency matrix are not defined. The definition of a node-neighbor
edge does not imply a direction of the communication.
\end{users} 

\item
Section~\ref{subsec:topol-inquiry} on page~\pageref{subsec:topol-inquiry}.
% --- MPI-2.1, Ballot 4, Item 10.b.i ---
\newline
In \mpifunc{MPI\_CARTDIM\_GET} and \mpifunc{MPI\_CART\_GET}: 
If \mpiarg{comm} is associated with a zero-dimensional Cartesian topology,
\mpifunc{MPI\_CARTDIM\_GET} returns \mpiarg{ndims=0} and \mpifunc{MPI\_CART\_GET} will keep
all output arguments unchanged. 

\item
Section~\ref{subsec:topol-inquiry} on page~\pageref{subsec:topol-inquiry}.
% --- MPI-2.1, Ballot 4, Item 10.b.ii ---
\newline
In \mpifunc{MPI\_CART\_RANK}:
If \mpiarg{comm} is associated with a zero-dimensional Cartesian topology,
\mpiarg{coord} is not significant and 0 is returned in \mpiarg{rank}. 

\item
Section~\ref{subsec:topol-inquiry} on page~\pageref{subsec:topol-inquiry}.
% --- MPI-2.1, Ballot 4, Item 10.b.iii ---
\newline
In \mpifunc{MPI\_CART\_COORDS}:
If \mpiarg{comm} is associated with a zero-dimensional Cartesian topology,
\mpiarg{coords} will be unchanged. 

\item
Section~\ref{subsec:topol-shift} on page~\pageref{subsec:topol-shift}.
% --- MPI-2.1, Ballot 4, Item 10.c ---
\newline
In \mpifunc{MPI\_CART\_SHIFT}:
It is erroneous to call \mpifunc{MPI\_CART\_SHIFT} with a direction that is either
negative or greater than or equal to the number of dimensions in the Cartesian
communicator. This implies that it is erroneous to call \mpifunc{MPI\_CART\_SHIFT} with a
\mpiarg{comm} that is associated with a zero-dimensional Cartesian topology.

\item
Section~\ref{subsec:topol-part} on page~\pageref{subsec:topol-part}.
% --- MPI-2.1, Ballot 4, Item 10.a ---
\newline
In \mpifunc{MPI\_CART\_SUB}: 
If all entries in \mpiarg{remain\_dims} are false or \mpiarg{comm} is already associated
with a zero-dimensional Cartesian topology then \mpiarg{newcomm} is associated
with a zero-dimensional Cartesian topology. 

% --- MPI-2.2 Ticket 61
\item[18.1.] 
Section~\ref{subsec:inquiry-version} on page~\pageref{subsec:inquiry-version}.
\newline
The subversion number changed from 0 to 1.
 
 
\item
Section~\ref{subsec:inquiry-inquiry} on page~\pageref{subsec:inquiry-inquiry}.
% --- MPI-2.1, Ballot 4, Item 13.i ---
\newline
In \mpifunc{MPI\_GET\_PROCESSOR\_NAME}: 
In C, a null character is additionally stored at \mpiarg{name[resultlen]}. 
\mpiarg{resultlen} cannot be larger then \mpiconst{MPI\_MAX\_PROCESSOR\_NAME}-1. 
In Fortran, name is padded on the right with blank characters. 
\mpiarg{resultlen} cannot be larger then \mpiconst{MPI\_MAX\_PROCESSOR\_NAME}. 

\item
Section~\ref{sec:errorhandler} on page~\pageref{sec:errorhandler}.
% --- MPI-2.1, Ballot 2, Item 3 ---
\newline
\mpifuncindex{MPI\_COMM\_GET\_ERRHANDLER}%
\mpifuncindex{MPI\_WIN\_GET\_ERRHANDLER}%
\mpifuncindex{MPI\_FILE\_GET\_ERRHANDLER}%
\mpiskipfunc{MPI\_\{COMM,WIN,FILE\}\_GET\_ERRHANDLER} behave as if a
new error handler object is created.
That is, once the error handler is no longer needed,
\mpifunc{MPI\_ERRHANDLER\_FREE} should be called with the error handler returned
from \mpifunc{MPI\_ERRHANDLER\_GET} or
\mpiskipfunc{MPI\_\{COMM,WIN,FILE\}\_GET\_ERRHANDLER}
to mark the error handler for deallocation.
This provides behavior similar to that of \mpifunc{MPI\_COMM\_GROUP} and
\mpifunc{MPI\_GROUP\_FREE}.

\item
Section~\ref{sec:inquiry-startup} on page~\pageref{sec:inquiry-startup}, see explanations to \mpifunc{MPI\_FINALIZE}.
% --- MPI-2.1, Ballot 1, Item 1 ---
\newline
\mpifunc{MPI\_FINALIZE} is collective over all connected processes.
If no processes were spawned, accepted or connected then this means
over \mpiconst{MPI\_COMM\_WORLD}; otherwise it is collective over the
union of all processes that have been and continue to be connected,
as explained in Section~\ref{subsec:disconnect} on page~\pageref{subsec:disconnect}.

\item
Section~\ref{sec:inquiry-startup} on page~\pageref{sec:inquiry-startup}.
% --- MPI-2.1, Ballot 4, Item 15.b ---
\newline
About \mpifunc{MPI\_ABORT}:
\begin{users} 
Whether the errorcode is returned from the executable or
from the \MPI/ process startup mechanism (e.g., mpiexec), is an aspect of
quality of the \MPI/ library but not mandatory.
\end{users} 
\begin{implementors} 
Where possible, a high-quality implementation will
try to return the errorcode from the \MPI/ process startup mechanism (e.g.
mpiexec or singleton init).
\end{implementors} 

\item
Section~\ref{subsec:info} on page~\pageref{subsec:info}.
% --- MPI-2.1, Ballot 4, Item 4 ---
\newline
An implementation must support info objects as caches for arbitrary (\mpishortarg{key},
\mpishortarg{value}) pairs, regardless of whether it recognizes the key. Each function that
\cdeclindex{MPI\_Info}%
takes hints in the form of an \mpictype{MPI\_Info} must be prepared to ignore any key it
does not recognize. This description of info objects does not attempt to
define how a particular function should react if it recognizes a key but not the
associated value. 
\mpifunc{MPI\_INFO\_GET\_NKEYS}, \mpifunc{MPI\_INFO\_GET\_NTHKEY},
\mpifunc{MPI\_INFO\_GET\_VALUELEN}, and \mpifunc{MPI\_INFO\_GET}
must retain all (\mpiarg{key},\mpiarg{value})
pairs so that layered functionality can also use the \mpiarg{Info} object. 

\item\label{item:changes:RMA:a}
Section~\ref{sec:onesided-putget} on page~\pageref{sec:onesided-putget}.
% --- MPI-2.1, Ballot 2, Item 6 ---
\newline
\mpiconst{MPI\_PROC\_NULL} is a valid target rank in the \MPI/ \RMA/ calls
\mpifunc{MPI\_ACCUMULATE}, \mpifunc{MPI\_GET}, and \mpifunc{MPI\_PUT}.  
The effect is the same as for \mpiconst{MPI\_PROC\_NULL} in \MPI/ point-to-point
communication.
See also item~\ref{item:changes:RMA:b} in this list. 

\item\label{item:changes:RMA:b}
Section~\ref{sec:onesided-putget} on page~\pageref{sec:onesided-putget}.
% --- MPI-2.1, Ballot 4, Item 6 ---
\newline
After any \RMA/ operation with rank \mpiconst{MPI\_PROC\_NULL}, it is still necessary to
finish the \RMA/ epoch with the synchronization method that started the epoch. 
See also item~\ref{item:changes:RMA:a} in this list. 

\item
Section~\ref{sec:1sided-accumulate} on page~\pageref{sec:1sided-accumulate}.
% --- MPI-2.1, Ballot 2, Item 7 ---
\newline
\mpiconst{MPI\_REPLACE}
in \mpifunc{MPI\_ACCUMULATE}, 
like the other predefined operations, 
is defined only for the predefined \MPI/ datatypes.

\item
Section~\ref{sec:io-info} on page~\pageref{sec:io-info}.
% --- MPI-2.1, Ballot 4, Item 18 ---
\newline
About \mpifunc{MPI\_FILE\_SET\_VIEW} and \mpifunc{MPI\_FILE\_SET\_INFO}:
When an info object that specifies a subset of valid hints is passed to
\mpifunc{MPI\_FILE\_SET\_VIEW} or \mpifunc{MPI\_FILE\_SET\_INFO}, 
there will be no effect on
previously set or defaulted hints that the \mpiarg{info} does not specify. 

\item
Section~\ref{sec:io-info} on page~\pageref{sec:io-info}.
% --- MPI-2.1, Ballot 4, Item 17 ---
\newline
About \mpifunc{MPI\_FILE\_GET\_INFO}:
If no hint exists 
for the file associated with \mpiarg{fh}, 
a handle to a newly created info object is returned that
contains no key/value pair. 

\item
Section~\ref{sec:io-view} on page~\pageref{sec:io-view}.
% --- MPI-2.1, Ballot 2, Item 10 ---
\newline
If a file does not have the mode \mpiconst{MPI\_MODE\_SEQUENTIAL},
then \mpiconst{MPI\_DISPLACEMENT\_CURRENT} is invalid as \mpiarg{disp} in 
\mpifunc{MPI\_FILE\_SET\_VIEW}.  

\item
Section~\ref{subsec:ext32} on page~\pageref{subsec:ext32}.
% --- MPI-2.1, Ballot 1, Item 7 ---
\newline
The bias of 16 byte doubles was defined with 10383. The correct value is 16383. 

\item
%-REMOVED-WITH-MPI-3.0: Section~\ref{sec:c++ClassMemberFunctions} on page~\pageref{sec:c++ClassMemberFunctions}.
\MPIIIDOTII/, Section 16.1.4 (Section was removed in \MPIIIIDOTO/).
% --- MPI-2.1, Ballot 1, Item 8 ---
\newline
In the example in this section,
the buffer should be declared as \code{const void* buf}. 

\item
Section~\ref{f90:types} on page~\pageref{f90:types}.
% --- MPI-2.1, Ballot 4, Item 16 ---
\newline
About \mpiskipfunc{MPI\_TYPE\_CREATE\_F90\_\XXX/}:
\begin{implementors}
An application may often repeat a call to
\mpiskipfunc{MPI\_TYPE\_CREATE\_F90\_\XXX/} with the same combination of 
(\mpiarg{\XXX/},\mpiarg{p},\mpiarg{r}).
The application is not allowed to free the returned predefined, unnamed
datatype handles. To prevent the creation of a potentially huge amount of
handles, the \MPI/ implementation should return the same datatype handle for
the same (\mpiarg{REAL/COMPLEX/INTEGER},\mpiarg{p},\mpiarg{r}) combination. 
Checking for the
combination (\mpiarg{p},\mpiarg{r}) in the preceding call to 
\mpiskipfunc{MPI\_TYPE\_CREATE\_F90\_\XXX/} and
using a hash-table to find formerly generated handles should limit the
overhead of finding a previously generated datatype with same combination
of (\mpiarg{\XXX/},\mpiarg{p},\mpiarg{r}).
\end{implementors}

\item
Section~\ref{subsec:annexa-const} on page~\pageref{subsec:annexa-const}.
% --- MPI-2.1, Ballot 3, Item 12 ---
\newline
\mpiconst{MPI\_BOTTOM} is defined as
\code{void * const MPI::BOTTOM}. 
 
\end{enumerate}
